\documentclass[
ngerman,
paper=a4,
asymmetric,
fontsize=10pt,
%headings=small,
draft=false,
% parskip=half,
numbers=noenddot,
% listof=totoc,         % verschiebt Minitoc/Hyperref
% bibliography=totoc,   % verschiebt Minitoc/Hyperref
bibtotocnumbered,
% index=totoc,
% noonelinecaption,
footnotes=multiple,
% toc=flat,
]{scrreprt}
\setlength{\parindent}{15pt}

\usepackage{luatextra}
\usepackage{fontspec}


\usepackage{graphicx} % Bilder-Support

%%%%%%%%%%%%%%%%%%%%%%%%%%%%%%%%%%%%%%%%%%%%%%%%%%%%%%%%%%%%%%%%%%%%%%%%%%%%%%%
%%%%%%%%% OPTIONEN
%%%%%%%%%%%%%%%%%%%%%%%%%%%%%%%%%%%%%%%%%%%%%%%%%%%%%%%%%%%%%%%%%%%%%%%%%%%%%%%
% optional --- Optionaler Text. 
\usepackage
[
%%% Status
% STATUS-DRAFT,
% STATUS-FINAL,
STATUS-EXPERIMENTAL,
% STATUS-DEV,
% STATUS-REVIEW,
% STATUS-FINAL,
%%% Git
%SHOW-GIT-COMMIT-ID,
%%%
% SHOW-TODO,
%%% Organisation
ORGASBFLO,
ORG-ASBFLO,
% ORG-ASBOE,
%%% Varianten
VAR-STANDARD,
VAR-RETTUNGSDIENST,
VAR-VAR0, % Stammvariante
% VAR-VAR1, % AASS, ASB Floridsdorf-Donaustadt
% VAR-VAR2, % ASBOE
%%% Region
REGAUT,
REG-AUT,
% REGGER;
REGVIE,
% REGNOE,   
]{optional}

\usepackage{microtype}
% \usepackage{fixltx2e}
% \usepackage{ellipsis}
% \setcounter{errorcontextlines}{10}

\hbadness=9000

% Sprachunterstützung 
\usepackage{polyglossia} % statt babel (LuaTeX)
	\setdefaultlanguage[
	spelling=new,
	latesthyphen=true,
	script=latin,
	]{german}
	\setotherlanguages{english}

% Mittlerweile verwendet pdfTeX standardmäßig PDF-1.5 mit
% Objektstreamkomprimierung, mit der pdftk Probleme haben kann.
% Füge \pdfobjcompresslevel=0 ein, dann wird die Datei zwar
% wieder größer, aber du kannst sie mit pdftk splitten.
% (Heiko Oberdiek)
\pdfobjcompresslevel=0

%%%%%%%%%%%%%%%%%%%%%%%%%%%%%%%%%%%%%%%%%%%%%%%%%%%%%%%%%%%
%%% Seitenlayout
\usepackage[
verbose,
a4paper,
%asymmetric,
% reversemp,
includemp,   % incl. marginpar
includehead, % incl. header
includefoot, % incl. footer
inner  = 30pt,
outer  = 30pt,
top    = 10mm,
bottom =  8mm,
bindingoffset=45.50787pt,%0.50787+42.67912+2 → Rest=552pt 
marginparsep=15pt,
marginparwidth = 154pt,
% marginparpush=0pt,
footskip = 7mm,
% % headheight=0.5cm,
% % headsep=0.5cm,	
% showframe,
]
{geometry}


% No bottom titles !!!
	\usepackage[nobottomtitles]{titlesec}
	\renewcommand{\bottomtitlespace}{.10\textheight}

\usepackage{hyperref}

%%% AASS Style Classes (ASC)
	\usepackage[AASS,stylelayout]{AassStyleCollectionCommon}
	\usepackage{AassStyleCollectionLayoutAlternative}
	\usepackage{AassStyleCollectionUnits}
	\usepackage{AassStyleCollectionFonts}
	\opt{VAR-VAR2}{\usepackage[MainFont=serif]{asboe}}
	\usepackage{AassStyleCollectionFeatures}
	\usepackage{AassStyleCollectionText}

%%% Spezielles Überschriftenformat
\usepackage{chngcntr}
\counterwithout{paragraph}{subsubsection} % removes paragraph from the 
%	\renewcommand{\thesection}		
%{{\color{DefaultA}\thechapter}{\color{DefaultB}.\arabic{section}}}
%	\renewcommand{\thesubsection}	
%{{\color{DefaultA}\thechapter}{\color{DefaultB}.\arabic{section}.\arabic{subsection}}}

%\counterwithin{paragraph}{chapter} % makes paragraph depend on chapter
\newcounter{TempParagraph}
\setcounter{TempParagraph}{1}

\DeclareRobustCommand{\TempParagraphUpdate}{%
\setcounter{TempParagraph}{\value{paragraph}}%
\addtocounter{TempParagraph}{1}%
} %\TempParagraphUpdate

	\renewcommand{\theparagraph}
	{{\color{AASSAnal3}\arabic{paragraph}.}}
%	\renewcommand{\thechapter}		
%	{\protect\TempParagraphUpdate\arabic{TempParagraph}.}
	\renewcommand{\thesection}		
	{\TempParagraphUpdate\color{AASSAnal3}\arabic{TempParagraph}.}
	\renewcommand{\thesubsection}	
	{\TempParagraphUpdate\color{AASSAnal3}\arabic{TempParagraph}.}
  \renewcommand{\thesubsubsection}	
  {\TempParagraphUpdate\color{AASSAnal3}\arabic{TempParagraph}.}

%subsubsections

%%% Abstand zwischen Absätzen in minipage-Umgebung:
% http://mrunix.de/forums/showthread.php?p=301707
	\makeatletter
	\newcommand*{\@minipagerestore}{\parskip@update}
	\makeatother






%%% Test PLAYGROUND


% obsolete? 
\newcommand{\VARIANT}[2]{{\opt{STATUS-DRAFT}{\color{SteelBlue}\sffamily\textbf{///~\uline{#1}$\rightarrow$}}\opt{#1}{#2}\opt{STATUS-DRAFT}{$\leftarrow$\textbf{\uline{#1}~///}}}}



% % Randnummernzähler einrichten 
% \newcounter{randnr} 
% \newcommand{\newrandnummer}{\refstepcounter{randnr}} 
% \newcommand{\rdnr}{\newrandnummer\marginnote{{\sffamily\bfseries\scriptsize\therandnr}}} 



%%% PaperTeX Adaption

	\newcommand{\papertexinexpandedtitle}[1]{
	\begin{minipage}{.95\textwidth}
	\begin{center}
	\noindent\Large\textbf{#1}
	\end{center}
	\end{minipage}
	}
	
	\newcommand{\expandedtitle}[2]{
	% \end{multicols}
	% \begin{center}
	\setlength{\fboxsep}{5pt}
	\setlength{\shadowsize}{2pt}
	\ifthenelse{\equal{#1}{shadowbox}}{%
	\shadowbox{%
	\papertexinexpandedtitle{#2}%
	}%
	}{}
	\ifthenelse{\equal{#1}{doublebox}}{%
	\doublebox{%
	\papertexinexpandedtitle{#2}%
	}%
	}{}
	\ifthenelse{\equal{#1}{ovalbox}}{%
	\ovalbox{%
	\papertexinexpandedtitle{#2}%
	}%
	}{}
	\ifthenelse{\equal{#1}{Ovalbox}}{%
	\Ovalbox{%
	\papertexinexpandedtitle{#2}%
	}%
	}{}
	\ifthenelse{\equal{#1}{lines}}{
	{\color{DefaultB}\hrule}
	\vspace*{1em}
	\begin{center}
	\noindent\sffamily\textit{#2}
	\end{center}
	\vspace*{1em}
	{\color{DefaultB}\hrule}
	}{}
	\ifthenelse{\equal{#1}{roundbox}}{
	% {\color{DefaultB}\hrule}
	% \vspace*{8pt}
	% \begin{center}
	% \\[0.5\baselineskip]\BoxRoundedBorderFilled{DefaultD}{White}{\center\hspace{0pt}\normalsize{\SlideHeading{#2}}}\\[0.5\baselineskip]
	\BoxRoundedBorderFilled{DefaultD}{White}{\center\hspace{0pt}\bfseries#2}
	% \vspace{-0.5\baselineskip}
	% \end{center}
	% \vspace*{8pt}
	%{\color{DefaultB}\hrule}
	}{}
	\ifthenelse{\equal{#1}{simple}}{
	\begin{center}
	\bfseries#2
	\end{center}
	}{}
	% \end{center}
	% \begin{multicols}{\papertex@ncolumns{}}
	% \ifnum \papertex@ncolumns > \minraggedcols
	% \raggedFormat
	% \fi
	}



\usetikzlibrary{shapes,calc}%Tbl of elements

\newcommand{\Study}[1]{\textit{#1}}

\definecolor{HelferA}{named}{AASSRed}
\definecolor{HelferB}{named}{AASSBlue}

% \newcommand{\License}[1]{
% 
% }

\opt{OPTIMIZATION-PRINT}{%
\settocfeature[toc][-1]{entryhook}{\FontSansNarrowFamily\normalsize\bfseries\color{BaseContrast}}
\settocfeature[toc][0]{entryhook}{\FontSansNarrowFamily\small\bfseries\color{Base}}
\settocfeature[toc][1]{entryhook}{\FontSansNarrowFamily\scriptsize}
\settocfeature[toc][2]{entryhook}{\FontSansNarrowFamily\scriptsize}
\settocfeature[toc][3]{entryhook}{\FontSansNarrowFamily\scriptsize}
}
\opt{OPTIMIZATION-ONLINE}{%
\settocfeature[toc][-1]{entryhook}{\sffamily\LARGE\bfseries\color{BaseContrast}}
\settocfeature[toc][0]{entryhook}{\sffamily\large\bfseries\color{Base}}
\settocfeature[toc][1]{entryhook}{\normalsize}
\settocfeature[toc][2]{entryhook}{\normalsize}
\settocfeature[toc][3]{entryhook}{\normalsize}
}


\makeindex

\input{./AASS.hyphenation.tex}
%%%%%%%%%%%%%%%%%%%%%%%%%%%%%%%%%%%%%%%%%%%%%%%%%%%%%%%%%%%
%%% Kopf- und Fußzeilen
\usepackage{fancyhdr}
\pagestyle{fancy}

\newcommand{\Fancyheadreset}{
  \fancyheadoffset[RE,RO]{\marginparsep+\marginparwidth}
  \fancyfootoffset[RE,RO]{\marginparsep+\marginparwidth}
  }


\pagestyle{fancy}
\fancyhf{} % clear all header and footer fields
\fancyheadoffset[RE,RO]{\marginparsep+\marginparwidth}
\fancyfootoffset[RE,RO]{\marginparsep+\marginparwidth}

\renewcommand{\chaptermark}[1]{\markboth{#1}{}}
\renewcommand{\sectionmark}[1]{\markright{\thesection.\ #1}}
\renewcommand{\headrule}{{\color{Grey}%
\hrule width\headwidth height\headrulewidth \vskip-\headrulewidth}}





%%% HEADER %%%
% right
\fancyhead[RO]{{\sffamily{\color{Grey}{\VarDocTitle} {\VarIterationNr}. Review und Iteration~|~}{\color{DefaultA}\thepage}}}
% \fancyhead[RE]{\EE{\protect\AuthNotes{ | {\protect\today} | ENTWURF |}}}
% left
% \fancyhead[LO]{\EE{\protect\AuthNotes{ | ENTWURF | {\protect\today} |}}}
\fancyhead[LE]{{\sffamily{\color{DefaultA}\thepage}{\color{Grey}~|~{\VarDocTitle} {\VarIterationNr}. Review und Iteration}}}
% %%% DRAFT %%%
\fancyhead[RE,LO]{{\noindent\sffamily\color{Grey}{\DocVersionInfo}}}

%%% FOOTER %%%
% center
\fancyfoot[CE,CO]{{{\sffamily\tiny  Zum internen Gebrauch}}}
% odd
\fancyfoot[LO]{{\noindent\sffamily\tiny\pdfdate}}
\fancyfoot[RO]{{\noindent{\sffamily\tiny\url{http://www.aass.at}~~|~~\VarDocTitle}~ \AASS{}}}
% even
\fancyfoot[LE]{{\noindent\AASS{}~ \sffamily\tiny \VarDocTitle~~|~~\url{http://www.aass.at}}}
\fancyfoot[RE]{{\noindent\sffamily\tiny\pdfdate}}






\makeindex

\usepackage[backend=bibtex8,
babel=other,
defernumbers=true,
natbib,
sorting=none,
sortlocale=de,
]{biblatex}

%   \bibliography{AASS.bib}
    \addbibresource{BiblioDB-Med-Books.bib}
    \addbibresource{BiblioDB-Med-Books-Harrisons18De.bib}
    \addbibresource{BiblioDB-Med-Journals-Rettungsdienst.bib}
    \addbibresource{BiblioDB-Med-Links.bib}
    \addbibresource{BiblioDB-Med-Misc.bib}
    \renewcommand{\labelnamepunct}{\addcolon\space}
    \DeclareFieldFormat{title}{\mkbibbold{\mkbibemph{#1}}}
    \DeclareFieldFormat[article]{title}{\mkbibbold{\mkbibemph{#1}}}
    \DeclareFieldFormat[inbook]{title}{\mkbibbold{\mkbibemph{#1}}}
    % \DeclareFieldFormat{citetitle}{\mkbibbold{#1}}
    % 
    % \DeclareFieldFormat{citetitle}[article]{\mkbibbold{#1}}
    \DeclareFieldFormat{journaltitle}{\textsc{#1}}
    \DeclareFieldFormat[article]{journaltitle}{\textsc{#1}}

%%%%%%%%%%%%%%%%%%%%%%%%%%%%%%%%%%%%%%%%%%%%%%%%%%%%%%%%%%%%%%%%   Begin Document

%\usepackage{cite}
%\bibliography{BiblioDB-Med-Misc.bib}
\usepackage{bibentry}
\nobibliography*

\newcommand{\Branch}[1]{\emph{#1}}


\usepackage[markup=underlined,ulem=normalem]{changes}
% \def\added{}
% \def\deleted{}
% \def\replaced{}

\newcommand{\Neu}{\fontfamily{lmtt}\fontseries{lc}\selectfont\color{DarkGreen}}
\newcommand{\ChangeHeading}[1]{\SlideHeading{#1}\addcontentsline{toc}{subsubsection}{#1}}

% \begin{itemize}
%   \item \TabHeading{Input}:
%   \begin{enumerate}
%     \item 
%   \end{enumerate}
%   \item \TabHeading{Verantwortlicher}:
%   \begin{enumerate}
%     \item 
%   \end{enumerate}
%   \item \TabHeading{Mitwirkende:}
%   \begin{enumerate}
%     \item 
%   \end{enumerate}
%   \item \TabHeading{Output}:
%   \begin{enumerate}
%     \item 
%   \end{enumerate}
% \end{itemize}

\title{Reviewprotokoll}
\subtitle{für die Arbeits- und Ausbildungsstandards für den
Sanitätsdienst}
\subject{44. Review und Iteration,\\ 2013-11-19}
\author{\AASS}


\newcommand{\VarDocTitle}{Vorläufiges Reviewprotokoll}
\renewcommand{\VarDocTitle}{Reviewprotokoll}
\newcommand{\VarIterationNr}{44}
\newcommand{\VarReviewDate}{2013-11-19}
\newcommand{\DocVersionInfo}{IN ARBEIT}
% \renewcommand{\DocVersionInfo}{FREIGEGEBEN}

\usepackage{datetime}


%\usepackage{lstdoc}
\usepackage{listings}
%\usepackage{textcomp}

% ============================
% EIGENE KOMMANDOS EMHOFER
% alles andere ist zu mühsam!!!
% ============================
\newcommand{\jALT}{\vspace{6pt}\noindent{\color{red}\sffamily\textbf{DELETED}\hrulefill\newline}}
\newcommand{\jALTende}{\newline{\color{red}\rule{\linewidth}{2pt}\newline}}
\newcommand{\jNEU}{\vspace{6pt}\noindent{\color{green}\sffamily\textbf{ADDED}\hrulefill\newline}}
\newcommand{\jNEUende}{\newline{\color{green}\rule{\linewidth}{3pt}\newline}}

\newcommand{\jganzebreite}{\textwidth+\marginparsep+\marginparwidth}


\newcommand{\RedFlag}{{\color{DefaultRed}\scriptsize{\VarFlag}\hspace{-0.5em}}}
\newcommand{\Dictum}[2][]{
  \bigskip
  \begin{flushright}
  \begin{minipage}{\ThirdNetWidth}\scriptsize
  {\ttfamily\sloppy\color{AltBlueDark}#2}\vspace{\baselineskip}
  
  {\FontSansNarrowFamily\color{AltBlueVeryLight}#1}
  \end{minipage}
  \end{flushright}
  \bigskip\noindent
  }

% %%%%%%%%%%%%%%%%%%%%%%%%%%%%%%%%%%%%%%%%%%%%%%%%%%%%%%%%%%%%%%%%%%%%%%%%%%%%%%%%%%%%%%%%%%%%%%%%%%%%%%#
% BEGIN DOCUMENT %%%%%%%%%%%%%%%%%%%%%%%%%%%%%%%%%%%%%%%%%%%%%%%%%%%%%%%%%%%%%%%%%%%%%%%%%%%%%%%%%%%%%%%%
% %%%%%%%%%%%%%%%%%%%%%%%%%%%%%%%%%%%%%%%%%%%%%%%%%%%%%%%%%%%%%%%%%%%%%%%%%%%%%%%%%%%%%%%%%%%%%%%%%%%%%%#

\newenvironment{Beispiel}
{\begin{tcolorbox}[title=Beispiel]

}{\end{tcolorbox}}

\setkomafont{labelinglabel}{\normalfont\bfseries}


\begin{document}
\MultiColListStart{toc}{2}
 
\lstset{% general command to set parameter(s)
language=[LaTeX]TeX,
%language=C,
extendedchars=true,
inputencoding=utf8,
breaklines=true;
basicstyle=\ttfamily\small, % print whole listing small
keywordstyle=\color{ubuntured}\bfseries\underbar,
% underlined bold black keywords
identifierstyle=\color{darkBlue}\bfseries, % nothing happens
commentstyle=\color{Grey}, % white comments
stringstyle=\ttfamily, % typewriter type for strings
texcsstyle=\color{ubuntured},
moretexcs={textcolor,itemhook},
emph={
AASS,
ARGE,
A,
aCh,
AuthorNotes,
Ca,
Cave,
Definition,
E,
EE,
ERROR,
Fazit,
H,
Layout,
Married,
MarriedNG,
MarriedLI,
Massnahmen,
marginpar,
Parallel,
T,
TODO,
Z,
ZFootnote},
emphstyle=\color{Terminus}\bfseries,
%directivestyle=color{Red},
showstringspaces=false,
%morecomment=[s][\color{SandyBrown}]{\{}{\}},
%moredelim=[s][\color{DarkViolet}]{\{}{\}},
%morecomment=[l][\color{Grey}]{\%}
}

\newgeometry{verbose,
a4paper,
asymmetric,
reversemarginpar,
rmargin=13mm,
lmargin=13mm,
top    = 20mm,
bottom = 15mm,
marginparwidth = 0mm,
footskip = 5mm,
bindingoffset=1.5cm}
\savegeometry{FullPage}
\newgeometry{verbose,
a4paper,
asymmetric,
reversemarginpar,
rmargin=13mm,
lmargin=13mm,
top    = 8mm,
bottom = 8mm,
marginparwidth = 0mm,
footskip = 5mm,
bindingoffset=1.5cm}
\savegeometry{WholePage}
\restoregeometry



\begin{titlepage}
\loadgeometry{FullPage}\thispagestyle{empty}
\sffamily
% {\BoxRoundedBorderFilled{DefaultA}{White}
% {{\sffamily\bfseries\Large\color{DefaultA}{\DocVersionTyp}} \par {\ttfamily\textbf{\DocVersion} Var. ASB\,921 --
% {\DocVersionInfo}}}}

% {\BoxRoundedBorderFilled{DefaultA}{White}
% {
% ~\vfill

\bigskip

\begin{center}
{{\Huge\scalefont{1.2}%{\color{DefaultGreyVeryLight}\dotfill}
\BoxRoundedBorderFilled{White}{DefaultGrey}{\centering\color{DefaultGrey}\Huge\VarDocTitle}
\bfseries

\vfill 

{\color{DefaultA}Arbeits- und
Ausbildungsstandards}

\bigskip

% {\color{Grey}\hrule}
% 
% \bigskip

{\color{DefaultB}für den Sanitätsdienst}} 

\bigskip

\bigskip
% 
{{\Large\color{Grey} {\VarIterationNr}. Review und Iteration}}
% \opt{VAR-VAR2}{{\Large\color{Grey} Angepasste Version für den Arbeiter-Samariter-Bund Österreichs}}

% \opt{VAR-VAR2}{
% \begin{minipage}[m]{{8cm}}
% \flushright{\Large\color{AsboeGreyVeryLight}\textit{Angepasste Version für den} \\\color{AsboeGreyLight} Arbeiter-Samariter-Bund Österreichs}
% \end{minipage}
% \hspace{1cm}
% \begin{minipage}[m]{3cm}
% \includegraphics[width=3cm]{./icons/asb}
% \end{minipage}
% 
% }


% \bigskip
% 
% {\color{Grey}\hrule}
% \opt{VAR-VAR2}{{
% % \begin{center}
% \bigskip
% \includegraphics[width=3cm]{./icons/asb}
% % \end{center}
% }}
% \bigskip
% 
% \Large{\AASS}
}
\end{center}



% \vspace{1cm}

% \bigskip

% }}

\vfill

\begin{center}
% \opt{VAR-VAR1}{{\includegraphics[height=12cm]{./images/pallinger-christoph-aass/Anton_32987_v2-AASS-0112mm}}}
% \opt{VAR-VAR2}{{\includegraphics[height=12cm]{./images/pallinger-christoph-aass/Anton_32987_v2-AASS-0112mm}}}
% \includegraphics[width=0.8\linewidth]{./images/pallinger-christoph-aass/Beatmungsmaske_32906_v2-AASS-0112mm}
\includegraphics[width=6cm]{./icons/under-construction}
\end{center}

\vfill


\begin{center}
% \BoxRoundedBorderFilled{DefaultD}{White}{
% \begin{center}
% {\Large\color{Grey}Herausgeber} \vspace{0.5em}
% \\{\Large\color{DefaultB}\bfseries Gabriel {$\cdot$} Koch {$\cdot$} Emhofer {$\cdot$} Motal {$\cdot$} Steuer}
% % \vspace{0.5em}
% % \\{\color{DefaultB}\itshape\normalsize{Mit Beiträgen von Lena
% % Hirtler, Christof Koller, Michael Withofner u.\,a.}}
% \end{center}

\bigskip

\BoxRoundedBorderFilled{DefaultGreyVeryLight}{DefaultGrey}{
\begin{multicols}{2}
% {Mit Beiträgen von }
% \vspace{0.5em}
% 
% {\color{DefaultB}Michael Auer\\
% Lena Hirtler\\
% Christof Koller\\ 
% Michael Withofner u.\,a.}
Review vom
\bigskip
\definecolor{LabelItemBoxStandard}{named}{DefaultGrey}
\renewcommand{\labelitemi}{\textcolor{LabelItemBoxStandard}{\SymbolLabelItemI}}

\VarReviewDate
\bigskip

\textbf{Zur internen Verwendung!}
\columnbreak

{Stand}
\vspace{0.5em}\\
{\textbf{\pdfdate}

({\DocVersionInfo})
}


% 
% \today
\end{multicols}
% }
}
\end{center}



\bigskip
% \date{November 2009}

% \opt{VAR-VAR1}{{
\begin{center}
{{\ARGE}}
\end{center}
% }}




\end{titlepage}
%\layout
\dominitoc
\addchap{Einleitung}

\begin{labeling}[:]{xxxxxxxxxxxxxxxxx}
  \item[Anwesende] ~
  \begin{itemize}
   % \item Johannes \caps{Steuer}
    \item Sebastian \caps{Gabriel}
    \item Josef \caps{Emhofer}
    \item Roman \caps{Koch}
    \item Michael \caps{Motal}
  \end{itemize}
  \item[Reviewkoordinator] Josef \caps{Emhofer}, Sebastian \caps{Gabriel}
  \item[Ort] ASB 921
  \item[Schwerpunkte] Releaseplanung 1.0
  \item[Repository] Git Version 0
  \item[Von Commit \#] db1d091d66e3462a7dcf7ca3c0428032ca99c3c9
  
  (17. Oktober 2013)
  \item[Bis Commit \#] c7d5975674e4396c7876498205c6086c9277276d
  
  (20. November 2013)
  \item[Anmerkungen] Es wurden nach dem Review noch Änderungen durchgeführt, diese werden extra ausgwiesen.
  \item[Anmerkungen] Textpassagen, die neu hinzugefügt wurden sind in der
  Umgebung
  
  \jNEU
  \E{Text, der hinzugefügt wurde}
  \jNEUende
  
  angegeben und Textpassagen, die entfernt wurden, sind durch die Umgebung
  
  \jALT
  \E{Text, der entfernt wurde}
  \jALTende
  
  gekennzeichnet.
  


\end{labeling}
\begin{WidePar}\small
\setcounter{tocdepth}{3}
\tableofcontents
\end{WidePar}


\chapter{Formale Änderungen im Buch (Layout)}
\begin{enumerate}
  \item Inhaltsverzeichnis in San Narrow-Font (siehe PDF)
  \item 100er-Stelle der Kapitelnummern steigen mit jedem Band
\end{enumerate}

\chapter{Inhaltliche und strukturelle Änderungen im Buch}

\section{Allgemeines}
\begin{enumerate}
  \item Dr. Withofner aus den Credits herausgenommen.
  \item Neues Kapitel: Fehlermanagement (advanced)
  \item Vorwort überarbeitet
  
  \jALT
  \noindent Liebe Leserinnen und Leser!
  
  Das vorliegende Werk ist das Ergebnis einer seit
mehr als vier Jahren andauernden Arbeit. 
Ziel dieses Werkes ist es, das Wissen und die Welt des Sanitätsdienstes abseits
bereits breitgetretener Pfade fundiert darzustellen, und
damit sowohl eine Unterlage für die Ausbildung, als auch ein
Hilfsmittel für die Praxis zu schaffen.  
Besonderen Wert haben wir dabei 
auf das Hintergrundwissen und relevante Randbereiche gelegt.
Denn in einer Zeit, in der das praktische Wissen innerhalb weniger Jahre bereits
wieder veraltet ist, bildet dieses Fundament die Chance für eine
langjährige Tätigkeit und erlaubt selbstständiges, regel\E{geleitetes} Arbeiten
(anstatt strikten, regel\E{gebundenen} Handelns). 

In diesem Sinne entstanden in Zusammenarbeit mit dem Ausbildungszentrum des ASB
Florisdorf-Donaustadt ausgehend von der Ausbildungsverordnung 
österreichischer Rettungssanitäter und den Erfordernissen der entsprechenden
Kurse die {\AASS}. Sie bilden die Grundlage der Curricula und konnten sich
schon, etappenweise und unter kontrollierten Bedingungen eingeführt, seit über
zwei Jahren als Lern- und Ausbildungsunterlage
bewähren.

Oft wurden wir
gefragt \Speech{\enquote{Wann wird Euer Buch fertig sein?}} 
Nun, das wird es wohl nie sein. Die {\AASS} sind darauf ausgelegt,
ständig weiterentwickelt, gewartet und verbessert zu werden. Das medizinische
und organisatorische Umfeld wandelt sich ständig, sodass man es sich gar nicht
leisten könnte, still zu stehen.
Mit dieser Erstausgabe möchten wir nun erstmalig an die
breitere Öffentlichkeit herantreten, und hoffen 
auf wertvolle Anregungen und konstruktive Kritik durch die
Leser- und Anwenderschaft.

\textbf{Dieses Werk ist keine \enquote{Dienstanweisung}!}~~

Eines muss vorweg klargestellt werden: Auch wenn es sehr in
Mode gekommen ist, Abläufe zu standardisieren und
festzuschreiben, so ist dies keine taugliche
Herangehensweise, um komplexe Probleme in der Medizin oder
im Sanitätsdienst zu bewältigen. Zu komplex sind die
Variablen, und zu gering bzw. zu unangepasst an konkrete
Situationen ist das gesicherte Wissen, auf dem
festgeschriebene Prozesse basieren. Der menschliche Faktor,
insbesonders die jeweilige Erfahrungen, ist und bleibt der
entscheidende Faktor jeglichen Handelns. Kein Lehrbuch oder
standardisierter Prozess kann das ändern, sofern das Ziel
ein \enquote{gutes} Ergebnis, und keine statistische Kennziffer
ist.

In diesem Sinne ist der Titel dieses Werkes nicht im Sinne eines
standardisierten Prozessmanagements zu verstehen, sondern soll einen
\E{Basisstandard} bezüglich medizinisch-sanitätsdienstlichem Wissen und
Kompetenz bilden, der individuell -- je nach Situation bzw. durch vom Helfer
andernorts erworbenen Wissens -- grundsätzlich verbesser- und erweiterbar ist. Die {\AASS}
geben Empfehlungen, es obliegt dem Helfer diese gemäß seiner 
fachlichen Erfahrung und Kenntnisse zu interpretieren, einzusetzen oder
anzupassen -- 

\Speech{salus aegroti suprema lex}\footnote{salus aegroti suprema lex \Lang{lat.}: \Speech{Das
Wohl des Kranken ist oberstes Gebot.}}.
  \jALTende
  
  \jNEU
  Das vorliegende Werk ist das Ergebnis einer seit dem Jahr 2008 
andauernden Arbeit. 
In Zusammenarbeit mit dem
Ausbildungszentrum des ASB Florisdorf-Donaustadt
enstanden ausgehend von der Ausbildungsverordnung
österreichischer Rettungssanitäter und den Erfordernissen
der entsprechenden Kurse die ersten Vrsionen der {\AASS}.
Sie bildeten die Grundlage der Curricula und
konnten sich schon, etappenweise und unter kontrollierten 
Bedingungen eingeführt, als Lern- und Ausbildungsunterlage bewähren.

Oft wurden wir dabei
gefragt \Speech{\enquote{Wann wird Euer Buch fertig sein?}} 
Nun, das wird es wohl nie sein. Die {\AASS} sind darauf ausgelegt,
ständig weiterentwickelt, gewartet und verbessert zu werden. Das medizinische
und organisatorische Umfeld wandelt sich ständig, sodass man es sich gar nicht
leisten könnte, still zu stehen.
Mit jeder Ausgabe hoffen wir
auf wertvolle Anregungen und konstruktive Kritik durch die
Leser- und Anwenderschaft.

\paragraph{Das Ziel des Projekts: \enquote*{Empowerment}}~~
% 
% \parpic[r][r]{\includegraphics[width=2cm]{./images/teaser/affe-sama-ambas-schein_00800.jpg}}
% 
% 
% 
% 
% 
\mbox{}\marginpar{\includegraphics[width=\linewidth]{./images/teaser/affe-sama-ambas-schein_00800.jpg}

\FontSansNarrowFamily\small Algorithmenaffe oder
Fachpersonal?}
Das zentrale Konzept der \T{Aufklärung} 
wird von Kant als die 
\Speech{\enquote{Befreiung von der selbst verschuldete
Unmündigkeit}} definiert. In diesem
Geiste entstanden und entstehen die
AASS: Sie sollen einerseits dem Sanitätspersonal, egal ob ärztlich oder
nicht-ärztlich, eine souveräne, überlegte und gleichzeitig
zielgerichtete Betreuung von Patienten und Notfallpatienten
ermöglichen und somit ein
Hilfsmittel für die Praxis zu schaffen. 
Andererseits war uns ist es aber
auch unser erklärtes Ziel, das Fachpersonal zum
eigenständigen Denken, Verstehen von Zusammenhängen und
Hintergründen und kritischen Hinterfragen zu befähigen:
Gelebtes \enquote*{Empowerment}. 

\paragraph{Geleitet, nicht gebunden}
Auch wenn es in
Mode gekommen ist, Abläufe zu standardisieren und
festzuschreiben, so ist dies keine taugliche
Herangehensweise, um komplexe Probleme in der Medizin oder
im Sanitätsdienst zu bewältigen. Zu komplex sind die
Variablen, und zu gering bzw. zu unangepasst an konkrete
Situationen ist das gesicherte Wissen, auf dem
festgeschriebene Prozesse basieren. Der menschliche Faktor,
insbesonders die jeweilige Erfahrungen, ist und bleibt der
entscheidende Faktor jeglichen Handelns. Kein Lehrbuch oder
standardisierter Prozess kann das ändern, sofern das Ziel
ein \enquote{gutes} Ergebnis, und nicht das Erreichen einer
statistischen Kennziffer ist.
Es ist Ausdruck der Würde
und Wertschätzung gegenüber den im Sanitätsdienst tätigen 
Personen, ihnen einen umfassenden
 Einblick in jene Welt zu gewähren, in der sie
arbeiten. 

\paragraph{Zeitloses Wissen}
Besonderen Wert haben wir dabei 
auf das Hintergrundwissen und relevante Randbereiche gelegt.
Denn in einer Zeit, in der das praktische Wissen innerhalb weniger Jahre bereits
wieder veraltet ist, bildet dieses Fundament die Chance für eine
langjährige Tätigkeit und erlaubt selbstständiges,
\EE{regel\E{geleitetes} Arbeiten} (anstatt strikten,
regel\E{gebundenen} Handelns).


So sind der Titel dieses Werkes und darin
enthaltene Handlungsempfehlungen nicht im Sinne eines
standardisierten, rigiden Prozessmanagements zu verstehen,
sondern sie sollen einen \E{Basisstandard} bezüglich
medizinisch-sanitätsdienstlicher Versorgung, Wissen und
Kompetenz bilden, der individuell -- je nach Situation 
bzw. durch vom Helfer erworbenen Wissens -- 
jederzeit verbesser- und erweiterbar ist. 
Die {\AASS} geben Empfehlungen, es obliegt dem Helfer diese gemäß seiner 
fachlichen Erfahrung und Kenntnisse zu interpretieren, einzusetzen oder
anzupassen -- 

\Speech{salus aegroti suprema lex}\footnote{salus aegroti suprema lex \Lang{lat.}: \Speech{Das
Wohl des Kranken ist oberstes Gebot.}}.
\jNEUende
\item Kapitelreihenfolge: [RSP] vor [HKL] (A vor B)
\item Umbenennung [RSP] $\to$ [ATM] Störungen der Atemwege und der Atmung
\item Lizenzen: Ergänzungen

     Sie dürfen folgende Seiten eigenständig gemäß den
     Bedingungen von \prettyref{topic:lizenz} nutzen:
     \begin{enumerate}
       \item Doppelseite Standardisierte Patientenversorung
      \item \deleted{\FullPrettyRef{tab:pam-weiterfuehrende}}
      \item \added{Übersicht
      Weiterführende
      Maßnahmen \ldots}
     \end{enumerate}
     
    [PD]
\deleted{\emph{Public domain}, \emph{"`Gemeinfrei"'}.}
\added{\emph{Public domain}, \emph{gemeinfrei}.}
 Die Gemeinfreiheit bezeichnet alle Werke, die keinem Urheberrecht mehr unterliegen
 oder ihm nie unterlegen haben. Das im angloamerikanischen Raum anzutreffende
 Public Domain ist ähnlich, aber nicht identisch mit der europäischen
 Gemeinfreiheit.
\end{enumerate}

\section{[ALS]}
 
 \Layout{2}{Neuerungen 2010}{%
 Die Empfehlung des
\enquote*{Pflicht-Zyklus vor AED} beim unbeobachteten
 Kreislaufstillstand ist gestrichen
 worden, kann aber weiterhin angewendet werden.
\deleted{Um die
HDM-Pausen zu reduzieren, soll während der Ladephase des
Defibrillators die HDM fortgeführt werden, d.\,h. nach der
Analyse wird weiter massiert, bis der Defi bereit zur}
\added{Während der Ladephase des
Defibrillators soll die Herzdruckmassage (\I{HDM})
fortgeführt werden um die HDM-Pausen zu reduzieren,
d.\,h. nach der Analyse wird weiter massiert, bis der Defi bereit zur}
 Schockabgabe ist.  Bei
 vielen AED-Geräten kann es jedoch zu
 \EE{Problemen} kommen, sie interpretieren die
 reanimationsbedingten Störungen der EKG-Ableitung als
 Kontaktverlust bzw. Bewegungsartefakt und brechen den
 Ladevorgang ab.\footnote{Davon betroffen sind derzeit z.\,B. die
 Geräte der Schiller Fred{\texttrademark}-Serie, der
 Defigard{\texttrademark} 1002 und der
 LifePak{\texttrademark} 500. Bezüglich anderer Geräte liegen
 der Redaktion derzeit noch keine Informationen vor. Es muss
 jedoch davon ausgegangen werden, dass auch andere Modelle
 ein derartiges Verhalten zeigen können. \E{(Angaben ohne
 Gewähr. Stand: 2011-01-19)}} Dadurch wird die Schockabgabe
 verhindert. Daher kann diese Methode nicht generell
 empfohlen werden, sondern darf nur durchgeführt werden, wenn
 der Hersteller des Gerätes dies freigegeben hat.
 \Commentary[HDM während Defi lädt / Doppelte Warnung]{Eine
 Warnung betreffend der HDM während der Ladephase des Defis
 findet sich bereits im \CO{[BLS]}-Teil
 (\prettyref{chp:BLS}). Aufgrund der Tragweite wurde sie im
 \CO{[ALS]}-Teil praktisch
 wortwörtlich wiederholt.}
 
 \Cave{Manche Defibrillatoren brechen den Ladevorgang ab,
 wenn währenddessen eine Herzdruckmassage durchgeführt wird.}
 
 \Fazit{Während des Ladevorganges des Defibrillators darf nur
 bei Verwendung von \E{manuellen} Geräten eine
 Herzdruckmassage durchgeführt werden.}
 }{}{}{}{}
 
\ldots

\jALT
 \Layout{2}{Endotracheale Intubation}
{Die frühe endotracheale Intubation (Einführen eines Tubuses
in die Luftröhre) verliert an Stellenwert. Sie soll nur
angewendet werden, wenn der Anwender hochqualifiziert
(\Speech{\enquote{high level of skill {\ldots}}}
\cite{Erc2010Sec4}) und sich sehr sicher ist
(\Speech{\enquote{{\ldots} and confidence}}
\cite{Erc2010Sec4}), bzw. regelmäßige, laufende Erfahrung
mit der Technik hat (\Speech{\enquote{and has regular,
ongoing experience with the technique.}}
\cite{Erc2010Sec4}).
Ohne ausreichendes Training und Erfahrung wird die
Komplikationsrate als unakzeptabel hoch angesehen. Das
Aufschieben der endotrachealen Intubation bis zur Wiederkehr
des spontanen Kreislaufs, wird in den ERC 2010-Leitlinien
als Alternative ausdrücklich genannt.
\cite{Nolan2001,
Nolan2008,
Wang2006,
Wang2006a,
Stiell2004,%There was no improvement in the rate of survival
% with the use of advanced life support in any subgroup.
Lecky2008,
Erc2010Sec4}
\Commentary[Airwaymanagement / Frühe Endotracheale Intubation (fETI)]
{Die endotracheale Intubation (ETI) ist grundsätzlich eine
sehr anspruchsvolle Tätigkeit, besonders im
außer-klinischen Bereich; es gibt viele Fehlerquellen
(Intubation des Ösophagus, Intubation eines Bronchus,
iatrogene Verletzung beim Intubationsvorgang, \ldots).
Zahlreiche Studien haben eine hohe Inzidenz von
Fehlintubationen beschrieben.

Nolan 2001 \cite{Nolan2001}
gibt eine umfassende
Übersicht bezüglich Erfolgsraten und Trainingsstrategien.
Schon damals waren die dargestellten Ergebnisse nicht
zufriedenstellend, gleichzeitig aber durchwachsen. Im
weiteren zeitlichen Verlauf wurde die ETI immer weiter
hinterfragt, z.\,B.
\cite{Wang2006,Wang2006a,Wang2008,Stiell2004,Stiell2008}.

In 2008 fassten Nolan und Soar \cite{Nolan2008} zusammen,
dass der Nutzen der ETI bei der präklinischen Reanimation
nicht durch Studien gezeigt werden konnte, und dass einige
sogar das Gegenteil erkennen ließen.
(\Speech{\enquote{Prehospital studies fail to show any
benefit from tracheal intubation during cardiopulmonary
resuscitation and many show harm.}}~\cite{Nolan2008}) Ein
ebenso 2008 veröffentlichtes Cochrane-Review von Lecky
et~al. \cite{Lecky2008} kam zu der gleichen Schlussfolgerung
mit der Betonung, dass entsprechende, qualitative
hochwertige Studien fehlen.
Dass die Probleme der ETI nicht auf nicht-ärztliches
Personal beschränkt ist, zeigten Timmermann et.~al. 2007
\cite{Timmermann2007}. In dieser in Deutschland
durchgeführten Studie wurde eine hohe Rate an
Fehlintubationen (ösophageal, bronchial) durch Notärzte
festgestellt.

Vor diesem Hintergrund bleibt festzuhalten, dass eine
unterlassene Intubation nicht mit einer unterlassenen
Hilfeleistung gleichzusetzen ist.}
}{}{}{}{}{}
 
 \subparagraph{Kontrolle der Tubuslage}
 Das korrekte Platzieren des Tubus ist einer der kritischten
 Punkte bei der endotrachealen Intubation. Typische Fehler
 sind die Intubation der Speiseröhre (\E{ösophageale}
 Intubation) oder die Intubation eines Hauptbronchus
 (\E{bronchiale} Intubation, wenn der Tubus zu weit
 vorgeschoben wird; es wird dann nur \E{ein} Lungenflügel belüftet).
 
 Grundsätzlich gibt es präklinisch drei Möglichkeiten, die
 korrekte Lage des Tubus zu überprüfen:
 \begin{enumerate}
   \item Visuelle Darstellung mittels Laryngoskop
   \item Auskultation von Magen und Lungenflügel
   \item Messung der \ce{CO2}-Abatmung mittels Kapnograhie
 \end{enumerate}
 
  Grundsätzlich sollen immer \EE{alle drei
 Möglichkeiten} angewendet werden.
 Es wird jedoch zunehmend ein größerer Wert auf die
 \EE{Kapnographie} gelegt, um die Tubuslage zu kontrollieren,
 die CPR-Qualität abzuschätzen und eine frühzeitige Erkennung
der Rückkehr des spontanen Kreislaufs zu ermöglichen.
Besonders geeignet sind dazu Geräte, welche die
\ce{CO2}-Abatmung als Kurve darstellen. Wenn nach 6
 Beatmungen \ce{CO2} nachweisbar ist, kann man davon
 ausgehen, dass der Tubus in den unteren Atemwege zu liegen
 gekommen ist. Die Kapnographie kann jedoch \E{keine} Aussage
 darüber treffen, ob die \E{Trachea oder ein einzelner
 Bronchus} intubiert wurde.
 
 Wenn kein Kapnograph mit Kurvendarstellung vorhanden ist,
 wird in den ERC~2010-Leitlinien die bevorzugte Verwendung
 supraglottischer Hilfsmittel
 angeregt.\ZFootnote{\Speech{\enquote{In the absence of a
 waveform capnograph it may be preferable to use a
 supraglottic airway device when advanced airway management
 is indicated.}}~\cite[p.1322]{Erc2010Sec4}.
 \Commentary[Airwaymanagement / Kapnographie]{Zur
 Problematik der Verfügbarkeit der Kapnographie siehe
 auch 
 \cite{Genzwuerker2007}
 (\fullcite{Genzwuerker2007}) 
 \cite{Timmermann2007}
 und darin enthaltene
 Referenzen.
 }
 }
.
\jALTende

\jNEU
 \Layout{2}{Endotracheale Intubation}
{\DefinitionInPara{Als
\T[Intubation!endotracheale]{endotracheale Intubation}
(\T{ETI}) wird das Einführen eines Tubus in die Trachea
bezeichnet.} Sie ist eine Maßnahme des invasiven
Atemwegsmanagements und wird als Goldstandard hinsichtlich
der Qualität der Atemwegssicherung angesehen, dies
berücksichtigt aber nicht die Sicherheit der Anwendung.

Der Tubus wird mithilfe eines Laryngoskops unter Darstellung
der Stimmritze unter Sicht in die Luftröhre eingeführt.
An den gängigen Tuben ist ein Cuff angebracht, welcher ein
Abdichten des Atemweges ermöglicht\footnote{Manche Tuben,
z.\,B. für Klinkinder haben keinen Cuff.}.
 
 \subparagraph{Kontrolle der Tubuslage}
 Das korrekte Platzieren des Tubus ist einer der kritischten
 Punkte bei der endotrachealen Intubation. Typische Fehler
 sind die Intubation der Speiseröhre (\E{ösophageale}
 Intubation) oder die Intubation eines Hauptbronchus
 (\E{bronchiale} Intubation, wenn der Tubus zu weit
 vorgeschoben wird; es wird dann nur \E{ein} Lungenflügel belüftet).
 
 Grundsätzlich gibt es präklinisch drei Möglichkeiten, die
 korrekte Lage des Tubus zu überprüfen:
 \begin{enumerate}
   \item Visuelle Darstellung mittels Laryngoskop
   \item Auskultation von Magen und Lungenflügel
   \item Messung der \ce{CO2}-Abatmung mittels Kapnograhie
 \end{enumerate}
 
  Grundsätzlich sollen immer \EE{alle drei
 Möglichkeiten} angewendet werden.
 Es wird jedoch zunehmend ein größerer Wert auf die
 \EE{Kapnographie} gelegt, um die Tubuslage zu kontrollieren,
 die CPR-Qualität abzuschätzen und eine frühzeitige Erkennung
der Rückkehr des spontanen Kreislaufs zu ermöglichen. Die
Kapnographie ist die verlässlichste Methode zur
Lagekontrolle \cite{Grmec2002}. Besonders
geeignet sind dazu Geräte, welche die \ce{CO2}-Abatmung als Kurve darstellen. Wenn nach 6
 Beatmungen \ce{CO2} nachweisbar ist, kann man davon
 ausgehen, dass der Tubus in den unteren Atemwege zu liegen
 gekommen ist. Die Kapnographie kann jedoch \E{keine} Aussage
 darüber treffen, ob die \E{Trachea oder ein einzelner
 Bronchus} intubiert wurde.
 
 Wenn kein Kapnograph mit Kurvendarstellung vorhanden ist,
 wird in den ERC~2010-Leitlinien die bevorzugte Verwendung
 supraglottischer Hilfsmittel
 angeregt.\ZFootnote{\Speech{\enquote{In the absence of a
 waveform capnograph it may be preferable to use a
 supraglottic airway device when advanced airway management
 is indicated.}}~\cite[p.1322]{Erc2010Sec4}.
 \Commentary[Airwaymanagement / Kapnographie]{Zur
 Problematik der Verfügbarkeit der Kapnographie siehe
 auch 
 \cite{Genzwuerker2007}
 (\fullcite{Genzwuerker2007}) 
 \cite{Timmermann2007}
 und darin enthaltene
 Referenzen.
 }
 }
}{}{}{}{}{}
.
\jNEUende

Streitfrage:


\jNEU
\Layout{2}
{Streitfrage: Endotracheale Intubation oder Alternative?}
{
Die endotracheale Intubation (ETI) ist grundsätzlich eine
sehr anspruchsvolle Tätigkeit.  Es gibt viele
Fehlerquellen, beispielhaft seien die Intubation des
Ösophagus oder nur eines Bronchus, Verletzung
beim Intubationsvorgang oder ein Bronchspasmus genannt.
Zahlreiche Studien haben eine hohe Inzidenz von
Fehlintubationen beschrieben.
Speziell im 
außerklinischen Bereich erhöht sich der
Schwierigkeitsfaktor aufgrund der Gegebenheiten
(Platzangebot, vorhandenes Material, nicht nüchterne bzw.
nicht vorbereitete Patienten, Verletzungen,
Begleiterkrankungen, \ldots) deutlich
\cite{Timmermann2006,Thierbach2004}.

In diversen Studien wurden die Erfolgs- und
Komplikationsraten der ETI untersucht, dabei zeigen sich
wenig zufriedenstellende bis alarmierende Ergebnisse %
\cite{Nolan2001,
% Prehospital and resuscitative airway care: should the gold standard be reassessed?
% In the context of prehospital care and resuscitation,
% tracheal intubation has been regarded as the standard in
% airway treatment. The evidence for this status is rather
% weak. It does not take into account the level of training
% and experience of the personnel attempting intubation, and
% whether they use neuromuscular blockers. In unskilled
% hands, attempted tracheal intubation is harmful;
% unrecognized esophageal intubation is disastrous. When
% healthcare providers lack adequate skills in tracheal
% intubation, alternative airway devices, such as the
% laryngeal mask airway or the Combitube, may be better
% options than a simple facemask. Healthcare personnel using
% any of these devices should be adequately trained and
% maintain frequent practice.
Wang2011,% Out-of-hospital airway management in the
% United States.In this study characterizing
% out-of-hospital airway management across the United
% States, we observed low out-of-hospital ETI success rates.
Cobas2009, % Prehospital intubations and mortality: a level
% 1 trauma center perspective.
% This prospective study showed a 31\% incidence of failed
% PHI in a large metropolitan trauma center. We found no
% difference in mortality between patients who were properly
% intubated and those who were not, supporting the use of
% bag-valve-mask as an adequate method of airway management
% for critically ill trauma patients in whom intubation
% cannot be achieved promptly in the prehospital setting.
Jones2004, % Emergency physician-verified out-of-hospital
% intubation: miss rates by paramedics. 
% The rate of unrecognized, misplaced out-of-hospital
% intubations in this urban, midwestern setting was 5.8\%
Wang2006, % Paramedic intubation errors: isolated events or
% symptoms of larger problems?
% In our study population, we found that errors occurred in
% 22 percent of intubation attempts, with a frequency of up
% to 40 percent in selected ambulance systems. These
% findings indicate frequent errors associated with this
% life-saving technique. These events might be emblematic of
% larger issues in the structure and delivery of
% out-of-hospital emergency care.
Wang2006a, % Out-of-hospital endotracheal intubation: where
% are we?
% These studies highlight our limited understanding of
% out-of-hospital endotracheal intubation and the need for
% new strategies to improve airway support in the
% out-of-hospital setting.
Wang2008,
Stiell2004,
Stiell2008,
}. Die wichtigsten Ergebnisse dieser Studien sind:
\begin{itemize}
  \item Die Gefahr von Fehlintubationen, insbesonders von
  ösophagealen Fehllagen, ist hoch.
  Diese werden oft spät oder gar nicht
  erkannt (Risiko der unerkannten
  Tubusfehllage bei Patienten mit außerklinischem
  Kreislaufstillstand: \Range{0,5}{17}\PC*%
  \cite{Grmec2002, % 0,5%: Comparison of three different
  % methods to confirm tracheal tube placement in emergency intubation.
  % The esophageal intubations were followed by successful
  % endotracheal intubations without complications. The
  % capnometry (sensitivity and specificity 100\%) and
  % capnography (sensitivity and specificity 100\%) were
  % better than auscultation (sensitivity 94\% and specificity
  % 83\%) in confirming endotracheal tube placement in
  % non-arrest patients ( p<0.05). Capnometry was highly
  % specific (100\%) but not sensitive (88\%) for correct
  % endotracheal intubation in patients with cardiopulmonary
  % arrest (capnometry versus auscultation and capnometry
  % versus capnography, p<0.05).Capnography is the most
  % reliable method to confirm endotracheal tube placement in
  % emergency conditions in the prehospital setting.
  Lyon2010, % 2,4%: Field intubation of cardiac
  % arrest patients: a dying art?
  % The aetiology of cardiac arrest was medical in 95.2\%,
  % traumatic in 3.9\% and unrecorded in 0.9\%. Prehospital
  % intubation was attempted in 628 patients. Prehospital
  % intubation was successful in 573 patients. A significant
  % complication (multiple attempts, displaced endotracheal
  % tube or oesophageal intubation) occurred in 55 (8.8\%)
  % patients. 165 (20.8\%) patients survived to hospital
  % admission, of whom 110 had undergone prehospital
  % intubation. 55 patients who did not undergo prehospital
  % tracheal intubation survived to hospital admission.The
  % optimal method of maintaining an airway and ventilating an
  % OHCA patient has yet to be established. Prehospital
  % tracheal intubation for OHCA is associated with
  % significant complications and may reduce survival. The use
  % of tracheal intubation as a routine intervention should be
  % reconsidered. Ambulance services should consider adopting
  % alternative strategies in airway management.
  Jones2004, % 6%: Emergency physician-verified
  % out-of-hospital intubation: miss rates by paramedics.
  Pelucio1997, % Out-of-hospital experience with the syringe
  % esophageal detector device.
  % In the remaining 168 intubations, the EDD correctly
  % identified 5 of 10 esophageal intubations (sensitivity
  % 50\%; 95\% CI: 19\%, 81\%).
  % The EDD correctly identified 156 of 158 tracheal
  % intubations (specificity 99\%; 95\% CI: 96\%, 100\%).In
  % this paramedic field study, the EDD demonstrated poor
  % sensitivity for esophageal intubations.
  Jemmett2003, % 9%: Unrecognized misplacement of
  % endotracheal tubes in a mixed urban to rural emergency
  % medical services setting.
  % Of the studied patients, 12\% (13 of 109) were found to
  % have misplaced endotracheal tubes. For the patients with
  % unrecognized improperly placed tubes, 9\% (10 of 109) were
  % in the esophagus, 2\% (2 of 109) were in the right main
  % stem, and 1\% (1 of 109) were above the cords. Paramedics
  % serving urban and suburban areas did not perform
  % significantly better (p < 0.05) than intermediate-level
  % providers serving areas that are more rural.
  % The incidence of unrecognized misplacement of endotracheal
  % tubes by EMS personnel may be higher than most previous
  % studies, making regular EMS evaluation and the
  % out-of-hospital use of devices to confirm placement
  % imperative. The authors were unable to show a difference
  % in misplacement rates based on provider experience or
  % level of training.
  Katz2001, % Misplaced endotracheal tubes by paramedics
  % in an urban emergency medical services system.
  % On arrival in the ED, 25\% (27/108) of patients were found
  % to have improperly placed endotracheal tubes. Of the
  % misplaced tubes, 67\% (18/27) were found to be in the
  % esophagus, whereas in 33\% (9/27), the tip of the tube was
  % found to be in the hypopharynx, above the vocal cords. %
  Erc2010Sek04De,
  }
  \item Die Erfolgsraten bei wenig routinierten Anwendern
  sind schlecht\footnote{\Speech{\enquote{Die
  Misserfolgsraten der Intubation betragen in
  außerklinischen, wenig ausgelasteten Systemen
  mit Anwendern, die selten Intubationen durchführen, bis
  zu 50\PC*}}\cite{Erc2010Sek04De}. [408, 409]}.
  \item Die Probleme der ETI  betreffen auch ärztliches
  Personal\footnote{Timmermann et.~al. 
  2007 \cite{Timmermann2007}: In dieser in Deutschland
  durchgeführten Studie wurde eine hohe Rate an
  Fehlintubationen (ösophageal, bronchial) durch Notärzte
  festgestellt.}. \cite{Timmermann2007,Grmec2002, % 0,5%:
  % Comparison of three different methods to confirm
  % tracheal tube placement in emergency intubation.
  }).
  \item Der Schulungsaufwand ist groß: Zum Erlernen der
  Tätigkeit werden wenigstens 100 ETI gefordert, zur
  Aufrechterhaltung mindestens 10 ETI pro
  Jahr \cite{Timmermann2012}. Damit verbunden wäre eine
  Vollzeit-Trainingsphase von etwa einem halben Jahr
  \cite{Bernhard2012}.
\end{itemize}

Erschwerend kommt hinzu, dass ein Vorteil für den Patienten
durch die ETI \E{nicht} nachgewiesen ist:
In 2008 fassten Nolan und Soar \cite{Nolan2008} zusammen, dass der Nutzen der ETI
bei der präklinischen Reanimation
nicht durch Studien gezeigt werden konnte, und dass einige
sogar das Gegenteil erkennen ließen
(\Speech{\enquote{Prehospital studies fail to show any
benefit from tracheal intubation during cardiopulmonary
resuscitation and many show harm.}}~\cite{Nolan2008}). Dies
deckt sich auch mit Beobachtungen aus anderen Studien \cite{
Stiell2004,%There was no improvement in the rate of survival
% with the use of advanced life support in any subgroup.
Cobas2009, % Prehospital intubations and mortality: a level
% 1 trauma center perspective. 
}.
Ein 2008 veröffentlichtes Cochrane-Review kam zu der
gleichen Schlussfolgerung, mit der Betonung, dass entsprechende, qualitative
hochwertige Studien fehlen \cite{Lecky2008}. 


Die frühe endotracheale Intubation verliert somit an
Stellenwert.
Sie soll nur angewendet werden, wenn der Anwender hochqualifiziert
und sich sehr sicher ist,
bzw. regelmäßige, laufende Erfahrung
mit der Technik hat 
(\Speech{\enquote{It should be used
only when trained personnel are available to carry out the procedure with a
high level of skill and confidence. }} \cite{Erc2010Sec4}).
Ohne ausreichendes Training und Erfahrung wird die
Komplikationsrate als unakzeptabel hoch angesehen. Das
Aufschieben der endotrachealen Intubation bis zur Wiederkehr
des spontanen Kreislaufs, wird in den ERC 2010-Leitlinien
als Alternative ausdrücklich genannt.
\cite{Nolan2001,
Nolan2008,
Lecky2008,
Erc2010Sec4}
Vor diesem Hintergrund bleibt festzuhalten, dass eine
unterlassene Intubation nicht -- wie oft behauptet -- mit
einer unterlassenen Hilfeleistung, sondern mit einer
unterlassenen Tradition gleichzusetzen ist.


Der traditionellen ETI steht eine stetige
Entwicklung an alternativen Airwaysystemen gegenüber, welche eine effektive und
technisch einfachere Sicherung des Atemweges ermöglichen
(supraglottsiche Atemwegshilfen).
Zwar müssen diese Systeme Abstriche hinsichtlich des
Aspirationsschutzes machen, wiegen dies jedoch durch ihre
einfach und sichere Anwendung auf %
\cite{Schalk2010, %
% Out-of-hospital airway management by paramedics and emergency physicians using laryngeal tubes.
% The LT-D/LTS-D represents a reliable tool for prehospital
% airway management in the hands of both paramedics and
% emergency physicians. It can be used as an initial tool to
% secure the airway until ETI is prepared, as a definitive
% airway by rescuers less experienced with ETI or as a
% rescue device when ETI has failed.
% The placement time was < or =45s (n=120), 46-90s (n=20)
% and >90s (n=7).
Wiese2009, % The use of the laryngeal tube disposable (LT-D)
% by paramedics during out-of-hospital resuscitation-an observational study concerning ERC guidelines 2005.
% The LT-D was used in 46\% of all cardiac arrest
% situations. Overall, the LT-D was successfully inserted in
% more than 90\% of all cases on first attempt. In 95\% of
% all cases, no problems concerning ventilation of the
% patient were described.As an alternative airway device
% recommended by the ERC in 2005, the LT-D may enable airway
% control rapidly and effectively. Additionally, by using
% the LT-D, a reduced "no-flow-time" and a better outcome
% may be possible.
}. Wenn die erforderliche Routine und Sicherheit für die
Durchführung einer ETI nicht gegeben ist und die Indikation
für ein invasives Atemwegsmanagement besteht, sollen
supraglottische Atemwegshilfen als \E{primärer} Zugang
eingesetzt werden \cite{Timmermann2012}.

}
{}
{}
{}
{}

\begin{Literary}
\fullcite{Timmermann2012}

\fullcite{Lecky2008}

\fullcite{Timmermann2007}

\fullcite{Erc2010Sek04De}
\end{Literary}
.
\jNEUende

\subsection{Larynxtubus}
 Der \T{Larynxtubus{\texttrademark}} wurde von der Firma VBM
 Medizintechnik entwickelt und zählt derzeit zu den
 populärsten Produkten zur supraglottischen
 Atemwegssicherung. Im Gegensatz zur endotrachealen
 Intubation kann er blind und ohne Hilfsmittel eingeführt
 werden, das Ende kommt bei korrekter Anwendung in der
 Speiseröhre zu liegen. Er besitzt \E{ein}
 Lumen\footnote{Als \I{Lumen} bezeichnet man denn
 Innendurchmesser von Röhren, Kanülen und Kathetern.} und
 zwei Cuffs\footnote{\I{Cuff}:
ein mit Luft gefüllter Ballon.}\deleted{, manche Modelle (LT-S) verfügen über
ein zweites Lumen, über -welches eine Magensonde zur Ableitung des Mageninhaltes
gelegt werden kann. Der obere Cuff kommt nach dem Einführen
im Rachen zu liegen, während der untere am Ende die
Speiseröhre abdichtet.}
\added{. Es gibt Modelle (LT-S),
welche über ein zweites Lumen verfügen, über welches eine
Magensonde zur Ableitung des Mageninhaltes gelegt werden
kann (Magendrainagekanal). Trotz korrekter Anwendung kann es
aufgrund von Leckagen zur Magenüberblähung kommen, daher
sollen unbedingt, wenn verfügbar, nur Modelle mit einem
Magendrainagekanal eingesetzt werden.
\cite{Dengler2011,Mann2013}}

\added{Der obere Cuff kommt nach dem Einführen im Rachen zu
liegen, während der untere am Ende die Speiseröhre abdichtet.}
 Dazwischen liegt die Öffnung des Lumens, nach dem Einführen
 in Höhe des Kehlkopfes. Da ober- und unterhalb der
 Abzweigung zur Luftröhre abgedichtet ist, kann die Luft über
 den Kehlkopf und den Kehldeckel in die Luftröhre strömen und
 die Gefahr der Aspiration von Mageninhalt kann deutlich
 reduziert werden. Es sind verschiedene Größen je nach
 Körperlänge erhältlich.
 \cite{Asai2005} % The laryngeal tube.
 
 Im direkten Vergleich mit der Beatmung mittels
 Beatmungsmaske
 sind bei einer korrekten Anwendung des Larynxtubus{\texttrademark} folgende
 Vorteile gegeben:
 
 \begin{itemize}
  \item \deleted{Durch den geblockten Ösophagus wird ein Überblähen verhindert}
  \item \added{Durch den geblockten Ösophagus soll das Überblähen
   des Magens sowie die Aspiration von Mageninhalt
  verhindert werden.}
   \item Die Anwendung erfordert keine höhere manuelle
   Geschicklichkeit als die Anwendung der Beatmungsmaske und
   -beutel.
   \item Zusätzliche Hilfsmittel zur Anwendung sind in der
   Regel nicht erforderlich, wodurch die Gefahr von
   Verletzungen durch solche reduziert ist.
 \end{itemize}
 
 \jALT
\begin{WidePar}
\begin{multicols}{3}
 \begin{enumerate}\CorrectionListsBox
  \item \EE{Notarzt} alarmieren
   \item \EE{Lagerung}: Rückenlage, Kopf in
   \EE{Schnüffelstellung} (Jackson Position, Kopf
   leicht erhöht (umgedrehte Nierentasse), nicht überstreckt,
   nicht gebeugt)
   \item Mund öffnen mit dem
   \enquote*{Kreuzgriff} (Daumen und Zeigefinger der linken
   Hand, der Daumen am Unterkiefer, der Zeigefinger am
   Oberkiefer (bzw. Zähne) öffnen den Mund. Achtung bei
   kariösen oder wackeligen Zähnen!)
   \item \EE{Einführen}: LT mit der Spitze
   entlang des harten Gaumens einführen, ev. Unterkiefer
   anheben
   \item Zahnreihe innerhalb der Markierungen
   \item
   \EE{Cuffen}, Volumen nach Farbcodierung
   \item Manuelles
   \EE{Fixieren}
   \item \EE{Beatmen}
   \item \EE{Erfolgskontrolle}: sehen,
   hören, \ce{CO2}-Messung, Sp\ce{O2} steigt, gleichbleibend
   \begin{enumerate}
     \item Sehen: Thoraxexkursionen
    \item Hören: Auskultation
    Kein Geräusch über dem
    Magen, Atemgeräusche über beiden Lungenflügeln.

    Achtung: Nicht seitengleiches Atemgeräusch beim
    Beatmen über den LT heißt: Verlegung eines Hauptbronchus
    (Fremdkörper) oder Pneumothorax, ist jedoch kein
    Hinweis für einen zu tief oder falsch liegenden Tubus
    (wie bei der Auskultation nach der endotrachealen
    Intubation))!
    \item Messung: sofortige 2-stellige et\ce{O2}-Messung,
     Sp\ce{O2} steigend oder gleichbleibend, Verfärbung des EasyCap
     II
   \end{enumerate}
   \item Wenn Beatmung gut möglich: \EE{Fixierung}
   \item Vorsichtiges Einführen einer \EE{Magensonde} und Entlasten des
   Magens (Magensonde erst einführen, wenn der Patient
   verlässlich beatmet ist und der LT fixiert ist (kein Muss!)
   \item Maximal 2 Versuche, max. 30\Sec* pro Versuch
   \item Absaugung von Sekret aus LT: Sekret aus dem LT  darf
   abgesaugt werden. Vorher die Absaugtiefe mit dem Abstand
   Mundwinkel-Ohrläppchen bestimmen
 \end{enumerate}
\end{multicols}
\end{WidePar}
 .
 \jALTende

\jNEU
 \begin{enumerate}\CorrectionListsBox
  \item Ggfs. \EE{Notarzt} alarmieren
   \item \EE{Lagerung}: Rückenlage, Kopf in
   \EE{Schnüffelstellung} (Jackson Position, Kopf
   leicht erhöht (umgedrehte Nierentasse), nicht überstreckt,
   nicht gebeugt)
   \item Mund öffnen mit dem
   \enquote*{Kreuzgriff} (Daumen und Zeigefinger der linken
   Hand, der Daumen am Unterkiefer, der Zeigefinger am
   Oberkiefer (bzw. Zähne) öffnen den Mund. Achtung bei
   kariösen oder wackeligen Zähnen!)
   \item \EE{Einführen}: LT mit der Spitze
   entlang des harten Gaumens einführen, ev. Unterkiefer
   anheben
   \item Zahnreihe innerhalb der Markierungen
   \item
   \EE{Cuffen}, Volumen nach Farbcodierung
   \item Manuelles
   \EE{Fixieren}
   \item \EE{Beatmen}
   \item \EE{Erfolgskontrolle}: sehen,
   hören, \ce{CO2}-Messung, Sp\ce{O2} steigt, gleichbleibend
   \begin{enumerate}
     \item Sehen: Thoraxexkursionen
   \item Hören: Auskultation\begin{itemize}
      \item Kein Geräusch über dem
    Magen
    \item Atemgeräusche über beiden Lungenflügeln
    \end{itemize}
    Achtung: Nicht-seitengleiches Atemgeräusch beim
    Beatmen über den LT bedeutet: Verlegung eines
    Hauptbronchus (Fremdkörper) oder Pneumothorax, es ist
    jedoch \E{kein} Hinweis für einen zu tief oder falsch
    liegenden Tubus (im Gegensatz zur endotrachealen
    Intubation)!
    \item Messung: sofortige 2-stellige et\ce{CO2}-Messung,
     Sp\ce{O2} steigend oder gleichbleibend, Verfärbung des EasyCap
     II
   \end{enumerate}
   \item Wenn Beatmung gut möglich: \EE{Fixierung}
   \item Vorsichtiges Einführen einer \EE{Magensonde} und Entlasten des
   Magens (Magensonde erst einführen, wenn der Patient
   verlässlich beatmet ist und der LT fixiert ist (kein Muss!)
   \item Maximal 2 Versuche, max. 30\Sec* pro Versuch
   \item Absaugung von Sekret aus LT: Sekret aus dem LT  darf
   abgesaugt werden. Vorher die Absaugtiefe mit dem Abstand
   Mundwinkel-Ohrläppchen bestimmen
 \end{enumerate}
.
\jNEUende


% ATM
\section{[ATM]}
Das Kapitel [RSP] Respiratorische Störungen wurde umbenannt in [ATM] Störungen
der Atemwege und der Atmung

\subsection{Blutungen aus dem Nasenraum}
\Layout{2}
{{\TxBeschreibung}}
{}
{}
{}
{}
{}

\section{[GYN]}
Umbenennung der Unterabschnitts \E{Uterusatonie} in \E{Blutungen aus der
Nachgeburtsperiode}.

\section{[HKL]}
\subsection{ACS}
\jALT

\Layout{0}
{\TxABCDE}
{~
\begin{ABCDEText}
%   \item[\ABCDE{1}] 
%   %%%
%   \item[\ABCDE{2}] 
  %%%
  \item[\ABCDE{3}]  \ldots
  %%%
  \item[\ABCDE{4}] \ldots
   
   Bei Diabetikern kann es aufgrund der Nervenschädigungen
   auch zu einem schmerzlosen Verlauf kommen, man spricht
   dabei auch vom \enquote*{\T[stummer Infarkt]{stummen
   Infarkt}}\index{Infarkt, stummer}.
  Sonst steht normalerweise der \E{stechende/brennende},
  nicht atemabhängige \II[Brustschmerz]{Brust-} oder
  \II{Oberbauchschmerz} zusammen mit \II{Atemnot} und
  \II{Todesangst} im Vordergrund. Der Brustschmerz kann in
  Richtung \E{Magen}, rechter oder linker \E{Hand},
  \E{Kiefer} oder auch \E{Rücken} \EE{ausstrahlen}.
   
   \EE{Todesangst und Vernichtungsgefühl} sind
   häufig.
   Bei Auftreten eines \EE{kardiogenen Schocks} können die
   entsprechenden Schockzeichen beobachtet werden: Blässe,
   Kaltschweißigkeit, Hypotonie, etc.
   (\FullPrettyRef{topic:schock-kardiogener}).
   %%%
   \item[\ABCDE{B}] \EE{Atemnot} kommt häufig vor und ist ein
   Leitsymptom.  Begleitend kann man oft eine Zyanose
   beobachten.
   %%%
  \item[\ABCDE{C}] Es können Zeichen eines \E{kardialen
  Schocks} vorliegen. 
   
   Der Puls kann unregelmäßig sein.
   %%%
  \item[\ABCDE{D}]  Der Blutdruck ist normal oder hypoton.
   %%%
  \item[\ABCDE{E}]  das EKG kann wichtige Hinweise über
   Ausdehnung, Alter und Lokalisation des Infarkts bringen.
   Um die Dynamik beurteilen zu können, soll schon möglichst
  früh ein 12-Kanal-EKG geschrieben werden. Die
  Beurteilung ist jedoch nur entsprechend ausgebildeten
  Personal vorbehalten.
 %   %%%
 %   \item[\ABCDE{\ldots}] 
   %%%
   \item[\ABCDE{=}]  Ein Patient mit einem  akuten
   Koronarsyndrom ist grundsätzlich als vital bedroht
   einzuschätzen. eine ärztliche Diagnose ist auch
   hinsichtlich der Auswahl des richtigen Transportziel
   (Herzkatheterlabor?) entscheidend.
 \end{ABCDEText}
 }
 {
 \begin{ABCDESlide}
 %   \item[\ABCDE{1}] 
 %   %%%
 %   \item[\ABCDE{2}] 
   %%%
  \item[\ABCDE{3}] 
  %%%
  \item[\ABCDE{4}] 
   %%%
%   \item[\ABCDE{A}] 
   %%%
  \item[\ABCDE{B}] 
   %%%
  \item[\ABCDE{C}] 
   %%%
  \item[\ABCDE{D}] 
   %%%
  \item[\ABCDE{E}] 
 %   %%%
 %   \item[\ABCDE{\ldots}] 
   %%%
  \item[\ABCDE{=}] 
 \end{ABCDESlide}
 }
 {}
 {}
 {}
  .
  \jALTende
  
  \jNEU
  
  \Layout{0}
{\TxABCDE}
{~
\begin{ABCDEText}
%   \item[\ABCDE{1}] 
%   %%%
%   \item[\ABCDE{2}] 
  %%%
  \item[\ABCDE{3}]  \ldots
  %%%
  \item[\ABCDE{4}] \ldots
     
   Bei Diabetikern kann es aufgrund der Nervenschädigungen
   auch zu einem schmerzlosen Verlauf kommen, man spricht
   dabei auch vom \enquote*{\T[stummer Infarkt]{stummen
   Infarkt}}\index{Infarkt, stummer}.
   
   \EE{Todesangst und Vernichtungsgefühl} sind
   häufig.
   Bei Auftreten eines \EE{kardiogenen Schocks} können die
   entsprechenden Schockzeichen beobachtet werden: Blässe,
   Kaltschweißigkeit, Hypotonie, etc.
   (\FullPrettyRef{topic:schock-kardiogener}).
   %%%
   \item[\ABCDE{B}] \EE{Atemnot} kommt häufig vor und ist ein
   Leitsymptom.  Begleitend kann man oft eine Zyanose
   beobachten.
   %%%
 \item[\ABCDE{C}] Der Puls kann unregelmäßig sein.   
  Der Blutdruck kann in jede Richtung
  verändert sein. Bedenke: Ein hypertensiver Notfall kann
  Auslöser eines \Abk{ACS} sein. Es
  können Zeichen eines \E{kardiogenen Schocks} vorliegen. 
   
%   \item[\ABCDE{D}]
   %%%
  \item[\ABCDE{E}]  Das EKG kann wichtige Hinweise über
   Ausdehnung, Alter und Lokalisation des Infarkts bringen.
   Um die Dynamik beurteilen zu können, soll schon möglichst
  früh ein 12-Kanal-EKG geschrieben werden um Veränderungen
  im Verlauf zu dokumentieren.
  Die Beurteilung ist jedoch nur entsprechend ausgebildeten
  und geübten Personal vorbehalten.
 %   %%%
 %   \item[\ABCDE{\ldots}] 
   %%%
   \item[\ABCDE{=}]  Ein Patient mit einem  akuten
   Koronarsyndrom ist grundsätzlich als vital bedroht
   einzuschätzen. eine ärztliche Diagnose ist auch
   hinsichtlich der Auswahl des richtigen Transportziel
   (Herzkatheterlabor?) entscheidend.
 \end{ABCDEText}
 }
 {
 \begin{ABCDESlide}
 %   \item[\ABCDE{1}] 
 %   %%%
 %   \item[\ABCDE{2}] 
   %%%
  \item[\ABCDE{3}]  Leidend, 
  evtl. kaltschweißig, blass
   %%%
  \item[\ABCDE{4}] \EE{Akuter,
  nicht atemabhängiger, Brustschmerz} (brennend, stechend;
  Hinter dem Brustbein, evtl.
   \EE{Ausstrahlung} in Kiefer, (linken) Arm, in Rücken)
  und
  \EE{Atemnot}, oft \EE{Enge-} bzw.
  \EE{Druckgefühl}, Todesangst
  
 

  Typisch untypisch:
  \EE{Oberbauchschmerz};  
  \enquote*{{stummer
  Infarkt}}  
   %%%
  \item[\ABCDE{B}] \EE{Atemnot}, evtl. Zyanose
   %%%
  \item[\ABCDE{C}] Puls evtl. unregelmäßig,   
  Blutdruck in jede Richtung
  verändert, evtl. Zeichen eines {kardiogenen Schocks}
  
   %%%
%   \item[\ABCDE{D}]
   %%%
  \item[\ABCDE{E}] EKG (im Verlauf)! 
 %   %%%
 %   \item[\ABCDE{\ldots}] 
   %%%
  \item[\ABCDE{=}]  Vital bedroht
 \end{ABCDESlide}
 }
 {}
 {}
 {}
  .
  \jNEUende
  
\subsection{Hypertensive Krise}
Änderung der Formulierung Buchstabe L beim SAMPLER-Schema:

\jALT
\ldots
\begin{SAMPLERText}
  \item[\SAMPLER{L}] Achte auf die letzte
  Medikamenteneinnahme sowie die letzten ärztlichen
  Kontrollen.
\end{SAMPLERText}
\ldots
\jALTende

\jNEU
\ldots
\begin{SAMPLERText}
   \item[\SAMPLER{L}] Wann war die letzte
  Medikamenteneinnahme bzw. die letzten ärztlichen
  Kontrollen? 
\end{SAMPLERText}
\ldots
\jNEUende


\section{[INK]}
Aufzählungen in der Randspalte ergänzt:

\subsection{Schweinegrippe, Vogelgrippe}
 \Layout{2}{Schweinegrippe, Vogelgrippe}
 {\Z{%
 Grippeviren befallen nicht nur den Menschen, auch
 (Wasser-)Vögel und andere Tiere erkranken an der Grippe.
 Normalerweise erfolgt keine Übertragung vom Tier zum
 Menschen. Überschreitet ein Erreger diese Grenze wird es
 gefährlich: Die Tierwelt beinhaltet dann einen nahezu
 unbegrenzten Vorrat an Keimen (Keimreservoir). Dieser kann 
 praktisch nicht kontrolliert werden, schließlich lassen
 sich Zugvögel nicht so leicht unter Quarantäne stellen
 \ldots 
 
 Glücklicherweise entwickelten sich in letzter Zeit nur
 Viren, welche zwar vom Tier auf den Menschen und dann
 weiter auf den Menschen übertragbar waren, jedoch waren sie
 (bis jetzt) nicht übermäßig infektiös oder zeigten keinen
 besonders schweren Verlauf. 
 
 \subparagraph{Die Gefahr} Die
 Tier\,\Sign{→}\,Mensch\,\Sign{→}\,Mensch-Übertragung stellt eine ständige
 Gefahr dar. 
 
 Ein weiteres Problem ist jedoch die mögliche Kreuzung mit
 Stämmen die sehr infektiös und schwer verlaufend sind. Dies
 kann passieren wenn ein Patient mit mehreren Stämmen
 gleichzeitig infiziert ist. Daher ist es wichtig die
 Infektionen so rasch als möglich einzudämmen.
 }
 
 \Fazit{Die große Gefahr ist die Entwicklung eines hoch
 infektiösen und schwer verlaufenden Stammes, welcher eine
 Tier-Mensch-Mensch-Übertragung aufweist.}
 \Fazit{Eine rasche Isolierung der Infizierten soll die
 Wahrscheinlichkeit einer bösartigen Kreuzung minimieren!}
 
 \subparagraph{Neue Grippe (\enquote{Schweinegrippe})} \Z{Die jüngste
 Entwicklung war die Übertragung des Stammes \Code{A(H1N1)}
 vom Schwein auf den Menschen in Mexiko im Frühjahr 2009. Der
 neue Stamm breitete sich sehr rasch um die ganze Welt aus,
 \E{die Erkrankung verlief jedoch recht mild}. Bemerkenswert
 war jedoch, dass -- anders als bei einer \enquote{normalen
 Grippe} üblich -- besonders junge Menschen infiziert wurden. 
 
 Man darf gespannt sein was die Zukunft bringt \ldots}
 }
{%
\jNEU
.
\Z{%
\begin{itemize}
  \item Keimreservoir: Tiere
  \item Gefahr: Kreuzung von Virenstämmen
  \item Gefahr: Mensch-Mensch-Übertragung
\end{itemize}
.
\jNEUende
}
 }
 {}
 {}
 {}

\subsection{Wundstarrkrampf}

 \Layout{20}{\TxSymptome}
 {%
 Es kommt zu tonischen \EE{Krämpfen der
 quergestreiften Muskulatur}, oft beginnend bei der
 Kaumuskulatur \Z{(Kiefersperre, Teufelsgrinsen
 (\I{Risus sardonicus}))}. Der Patient hat starke Schmerzen
 und entwickelt in Folge der Krämpfe eine \EE{Atemlähmung}. 
 Die Fortschritte der
 Intensivmedizin haben die Behandlungsmöglichkeiten
 verbessert, trotzdem sterben noch etwa \VonBis{20}{30}\PC*
 der Tetanuserkrankten. Der Wert der
 Impfung ist daher nach wie vor groß.
 }
 {
 \jNEU
 .
\begin{itemize}
  \item Krämpfe der Skelettmuskulatur
  \item Schmerzen, Atemlähmung
  \item Hohe Letalität
\end{itemize}
.
\jNEUende
 }
 {./images/hygiene/6373_lores_tetanus.jpg}
 {Patient mit Tetanus.}
 {Public Health Image Library, ID\#: 6373, Public Domain.}

\section{[KIN]}
Umbenennung des Unterabschnitts \E{Laryngitis, Pseudokrupp} in \E{Akute
obstruktive Laryngitis, Pseudokrupp}.

\section{[NEU]}
\subsection{Insult: Spezielle Anamnese}
\jNEU
\Layout{0}{Spezielle Anamnese}
{%
Vor Anmeldung zu einer Stroke-Unit müssen folgende
Informationen erhoben worden sein:

\begin{itemize}
  \item Zeitpunkt des Auftretens der Symptome
  \item 
\end{itemize}

}{
\begin{itemize}
  \item 
\end{itemize}

}{}{}{}
.
\jNEUende



\section{[PAM]}
\subsection{Einschätzungsblock:
1\,--\,2\,--\,3\,--\,4\,--\,A\,--\,B\,--\,C\,--\,D\,--\,E}

\jNEU
\begin{quote}
  \Speech{\enquote{Von 1 bis E.}}
\end{quote}
.
\jNEUende

\ldots

 \subparagraph{Aber was ist mit {\ldots}?}
Das ABCDE-Schema bzw. der Einschätzungsblock
\added{\T{\Range{1}{E}}} dient der Erkennung von vital bedrohten 
Patienten und dem Erkennen von Hauptbeschwerden und
 Leitsymptomen, \E{es deckt nicht sämtliche
 mögliche Untersuchungen ab!} Weiterführende Untersuchungen
 und Anamneseschritte bzw. Maßnahmen werden \E{nach} dem
 Einschätzungsblock nach Bedarf durchgeführt.


\section{[TEA]}
Aufzählungen in Randspalten hinzu gefügt.

\Layout{0}{Theoretische Anforderungen}{% 
 Die Tätigkeit im medizinischen Bereich geht
 aufgrund der Komplexität der Aufgaben mit einer Fülle von
 körperlichen, wissensmäßigen und psychischen Anforderungen
 einher. Einen nur kargen Hinweis darüber, welche Qualitäten
 z.\,B. von einem Sanitäter erwartet werden, liefert das
 Anforderungsprofil des Arbeitsmarktservice (AMS):
 \begin{itemize}
 \item Körperliche und psychische Belastbarkeit
 \item rasches Auffassungs- und Reaktionsvermögen
 \item Einfühlungsvermögen
 \item Beobachtungsgabe
 \item Kommunikations- und Teamfähigkeit
 \item Zuverlässigkeit
 \item Verantwortungsbewusstsein
 \end{itemize}
 }{%
\jNEU
.
\begin{itemize}
\item Körperliche und psychische Belastbarkeit
\item rasches Auffassungs- und Reaktionsvermögen
\item Einfühlungsvermögen
\item Beobachtungsgabe
\item Kommunikations- und Teamfähigkeit
\item Zuverlässigkeit
\item Verantwortungsbewusstsein
\end{itemize}
.
\jNEUende
 }{}{}{}

\subsection{Ärzte}
\jALT
Der ärztliche Beruf wird durch das \I{Ärztegesetz} (\I{ÄrzteG}) geregelt.
 Der Arzt ist zur Ausübung der Medizin berufen. Dies umfaßt
 jede auf medizinisch-wissenschaftlichen Erkenntnissen
 begründete Tätigkeit, die unmittelbar am Menschen oder
mittelbar für den Menschen ausgeführt wird, insbesondere 
die Untersuchung auf das Vorliegen oder Nichtvorliegen von körperlichen und psychischen Krankheiten oder Störungen,
die Beurteilung von diesen Zuständen bei Verwendung medizinisch-diagnostischer Hilfsmittel sowie deren Behandlung;
 die Vornahme operativer Eingriffe einschließlich der Entnahme oder Infusion von Blut;
 die Vorbeugung von Erkrankungen;
 die Geburtshilfe;
 die Verordnung von Heilmitteln, sowie
 die Vornahme von Leichenöffnungen.
 Ferner ist
der Arzt befugt ärztlichen Zeugnis auszustellen ärztliche
Gutachten zu erstatten.

Die selbstständige Ausübung des ärztlichen Berufes
ist ausschließlich Ärzten für Allgemeinmedizin und
approbierten Ärzten sowie Fachärzten vorbehalten
(Vorbehaltsregelung, von diesem Vorbehalt gibt es
für andere medizinische Berufsgruppen
entsprechende Ausnahmeregelungen).

Die in Ausbildung zum Arzt für Allgemeinmedizin oder zum
Facharzt befindlichen Ärzte (\I[Turnusarzt]{Turnusärzte}) sind lediglich zur
unselbstständigen Ausübung  in  als
Ausbildungsstätten anerkannten Einrichtungen, im Rahmen von
Lehrpraxen bzw. Lehrgruppenpraxen oder in Lehrambulatorien
unter Anleitung und Aufsicht der ausbildenden Ärzte
berechtigt.
 
\subsection{Ausbildung}

Die ärztliche Ausbildung gründet sich auf einem
 abgeschlossenes Universitätsstudium, der weitere
 Ausbildungsweg ist jedoch international sehr
unterschiedlich.
 
\subsubsection{Situation in Österreich}
 
In Österreich dauert das Studium der Humanmedizin je nach
 Universität und Curriculum \Range{5}{6} Jahre. Im Anschluss an das
 Studium verfügen die Absolventen \E{nicht} über das Recht zur
 Berufsausübung, sie sind auch nach dem
 Studium noch nicht als Ärzte, sondern als Mediziner zu
bezeichnen. Nach Abschluss des Studiums wird ein
 Ausbildungsverhältnis angetreten, es erfolgt die Aufnahme in
die Ärzteliste, welche von der Österreichischen Ärztekammer
geführt wird, der Arzt in Ausbildung darf nun an einer
anerkannten Ausbildungseinrichtungen unter 
 Aufsicht ärztlich tätig werden. Dies kann entweder im Rahmen
 eines Ausbildungsverhältnisses zum Arzt für Allgemeinmedizin,
 welches mindestens drei Jahre dauert, erfolgen oder in Ausbildung zu
 einem Facharzt, diese Ausbildungszeit beträgt je nach Fach
 \Range{5}{6} Jahre. Die Ausbildung zum Facharzt kann
 grundsätzlich gleich nach dem Studium beginnen, bzw. auch erst nach der
 Absolvierung der Ausbildung zum Arzt für Allgemeinmedizin
 begonnen werden. Während der Ausbildungszeit durchläuft der
 Auszubildende  im Turnus-System diverse
 Fachabteilungen, welche in der jeweiligen Ausbildungsordnung
vorgegeben sind (daher die Bezeichnung Turnusarzt). 
 
 Ein Arzt für Allgemeinmedizin verfügt über ein so genanntes
uneingeschränktes \I{Jus practicandi}, d.\,h. das Recht zu
praktizieren. Ein Facharzt verfügt ebenfalls über ein Jus
practicandi, welches jedoch oft Tätigkeiten in seinem
 Fachgebiet beschränkt ist, es sei denn, der betreffende
 Facharzt hat auch die Ausbildung zum Arzt für
 Allgemeinmedizin absolviert. Gegenwärtig gibt es in
 Österreich über 40 Fachgebiete für die man sich
qualifizieren kann.
 
Interessensvertretung und Aufsichtsbehörde für alle Ärzte
ist die \I{Österreichische Ärztekammer}\index{Ärztekammer} bzw. die jeweils
zuständigen Landesärztekammern.
 
Der Notarzt
 
Die Ausbildung zum Notarzt setzt die Absolvierung eines 60-stündigen, von der
österreichische Ärztekammer anerkannten und mit einer
Prüfung abschließenden Diplomkurs
voraus.
 
 
Zum Erhalt der
Tätigkeitsberechtigung ist eine regelmäßige
Rezertifizierung (Notarzt-Refresher-Kurs) notwendig.
 
Leitender Notarzt
 
Das Ärztegesetz sieht für leitende Funktionen, insbesondere
auch im Großschadensfall, die Ausbildung des leitenden
Notarztes vor. Die Ausbildung zum Leitenden Notarzt setzt
eine entsprechende Berufserfahrung als Notarzt sowie die
Absolvierung eines entsprechenden Kurs voraus.
\jALTende

\jNEU
 Der Arzt ist zur Ausübung der Medizin berufen. Dies umfaßt
 jede auf medizinisch-wissenschaftlichen Erkenntnissen
 begründete Tätigkeit, die unmittelbar am Menschen oder
mittelbar für den Menschen ausgeführt wird.\ZFootnote{Die
Ausübung der Medizin umfasst insbesonders die
Untersuchung auf das Vorliegen oder Nichtvorliegen von körperlichen und psychischen Krankheiten oder Störungen, die Beurteilung von diesen Zuständen bei Verwendung medizinisch-diagnostischer Hilfsmittel sowie deren Behandlung;
 die Vornahme operativer Eingriffe einschließlich der Entnahme oder Infusion von Blut;
 die Vorbeugung von Erkrankungen;
 die Geburtshilfe;
 die Verordnung von Heilmitteln, sowie
 die Vornahme von Leichenöffnungen.
 Ferner ist
der Arzt befugt, ärztliche Zeugnisse auszustellen und
ärztliche Gutachten zu erstatten.}
 
\Layout{2}
{Ausbildung}
{Die ärztliche Ausbildung gründet sich auf einem
 abgeschlossenes Universitätsstudium, der weitere
 Ausbildungsweg ist jedoch international sehr
unterschiedlich.}
{}
{}
{}
{}
 
\subsection{Situation in Österreich}

\Layout{2}
{Allgemeines}
{Der ärztliche Beruf wird durch das \I{Ärztegesetz} (\I{ÄrzteG}) geregelt.
Die selbstständige Ausübung des ärztlichen Berufes
ist ausschließlich \I[Arzt!für Allgemeinmedizin]{Ärzten für
Allgemeinmedizin}, \I[Arzt!approbierter]{approbierten
Ärzten} sowie \I[Facharzt]{Fachärzten}\index{Arzt!Fach-}
vorbehalten (\I{Arztvorbehalt}, von diesem Vorbehalt gibt es für andere medizinische Berufsgruppen
entsprechende Ausnahmeregelungen in den jeweiligen
Berufsgesetzen). \ZFootnote{Die in Ausbildung zum Arzt für
Allgemeinmedizin oder zum Facharzt befindlichen Ärzte (\I[Turnusarzt]{Turnusärzte}) sind lediglich zur
unselbstständigen Ausübung  in  als
Ausbildungsstätten anerkannten Einrichtungen, im Rahmen von
Lehrpraxen bzw. Lehrgruppenpraxen oder in Lehrambulatorien
unter Anleitung und Aufsicht der ausbildenden Ärzte
berechtigt.}

Interessensvertretung und Aufsichtsbehörde für alle Ärzte
ist die \I{Österreichische Ärztekammer}\index{Ärztekammer} bzw. die jeweils
zuständigen Landesärztekammern.}
{}
{}
{}
{}

\Layout{2}
{Ausbildung und Jus practicandi}
{In Österreich dauert das Studium der Humanmedizin je nach
 Universität und Curriculum \Range{5}{6} Jahre. Im Anschluss an das
 Studium verfügen die Absolventen \E{nicht} über das Recht zur
 Berufsausübung, sie sind auch nach dem
 Studium noch nicht als Ärzte, sondern als Mediziner zu
bezeichnen. 

Nach Abschluss des Studiums wird ein
 Ausbildungsverhältnis angetreten, es erfolgt die Aufnahme in
die \EE{Ärzteliste}, welche von der Österreichischen
Ärztekammer geführt wird. Der Arzt in Ausbildung darf nun an
einer anerkannten Ausbildungseinrichtungen unter 
 Aufsicht ärztlich tätig werden. Dies kann entweder im Rahmen
 eines Ausbildungsverhältnisses zum Arzt für Allgemeinmedizin,
 welches mindestens drei Jahre dauert, erfolgen oder in Ausbildung zu
 einem Facharzt, diese Ausbildungszeit beträgt je nach Fach
 \Range{5}{6} Jahre. Die Ausbildung zum Facharzt kann
 grundsätzlich gleich nach dem Studium beginnen, bzw. auch erst nach der
 Absolvierung der Ausbildung zum Arzt für Allgemeinmedizin
 begonnen werden. Während der Ausbildungszeit durchläuft der
 Auszubildende  im Turnus-System diverse
 Fachabteilungen, welche in der jeweiligen Ausbildungsordnung
vorgegeben sind (daher die Bezeichnung
\T{Turnusarzt}\index{Arzt!Turnus-}). Die Turnus-Ausbildung
zum Arzt für Allgemeinmedizin dauert 3 Jahre bzw. zum
Facharzt je nach Fachgebiet \Range{5}{6} Jahre. 
 
 Ein Arzt für Allgemeinmedizin verfügt über ein so genanntes
uneingeschränktes \T{Jus practicandi} (Recht zu
praktizieren). Ein Facharzt verfügt ebenfalls über ein Jus
practicandi, welches jedoch auf Tätigkeiten in seinem
 Fachgebiet beschränkt ist, es sei denn, der betreffende
 Facharzt hat auch die Ausbildung zum Arzt für
 Allgemeinmedizin absolviert. Gegenwärtig gibt es in
 Österreich über 40 Fachgebiete für die man sich
qualifizieren kann.}
{}
{}
{}
{}
 
\Layout{2}
{Der Arzt im Rettungs- und Krankentransportdienst}
{Das ÄrzteG sieht hierfür zwei Ausbildungen vor:
\begin{itemize}
  \item \T{Notarzt}: Die Ausbildung zum Notarzt setzt die
  Absolvierung eines 60-stündigen, von der österreichische Ärztekammer anerkannten und mit einer
  Prüfung abschließenden Diplomkurs
  voraus.
  Zum Erhalt der
  Tätigkeitsberechtigung ist eine regelmäßige
  Rezertifizierung (Notarzt-Refresher-Kurs) notwendig.
  \item \T{Leitender Notarzt}: Leitende Funktion, insbesonders 
  im Großschadensfall. Die Ausbildung zum Leitenden Notarzt setzt
  eine entsprechende Berufserfahrung als Notarzt sowie die   Absolvierung eines
  entsprechenden Kurs voraus.
\end{itemize}

}
{}
{}
{}
{}
.
\jNEUende

 \subsection{Gehobener Gesundheits- und Krankenpflegedienst}
 
 \jALT
 Zum gehobenen Gesundheits- und Krankenpflegedienst gehören
die diplomierten Gesundheits- und Krankenschwestern, die
diplomierten Gesundheits- und Krankenpfleger, sowie die
diplomierten Kinderkrankenschwestern und -pfleger (\I{DGKS},
\I{DGKP}, \I{DGKKS}, \I{DGKKP}). 
% TODO Psychiatriepflege
 Die Ausbildung erfolgt
an Fachschulen (\enquote*{Schwesternschule}), neuerdings
 wird die Ausbildung auch im Rahmen eines Fachhochschulstudiums angeboten. 
 Der Tätigkeitsbereich des gehobenen Gesundheits- und
 Krankenpflegedienstes beinhaltet einen
 eigenverantwortlichen 
Tätigkeitsbereich, welcher insbesond die Pflege von
 von Patienten vorsieht, sowie einen mitverantwortlich en
 Tätigkeitsbereich, welcher die Durchführung von einfachen
 ärztlichen Tätigkeiten (Blutentnahme, Verabreichung von
Infusionen und Medikamenetn, \ldots) vorsieht.
\jALTende

\jNEU
 Zum gehobenen Gesundheits- und Krankenpflegedienst gehören
die diplomierten Gesundheits- und Krankenschwestern und
-pfleger (\I{DGKS},
\I{DGKP}),
die diplomierten Kinderkrankenschwestern und -pfleger 
(\I{DGKKS}, \I{DGKKP}), sowie die diplomierten
psychiatrischen Gesundheits- und
Krankenschwestern und -pfleger.
 Die Ausbildung erfolgt
an Fachschulen (\enquote*{Schwesternschule}), aktuell
 wird die Ausbildung auch im Rahmen eines Fachhochschulstudiums angeboten. 
 Der Tätigkeitsbereich des gehobenen Gesundheits- und
 Krankenpflegedienstes beinhaltet einen
 eigenverantwortlichen 
Tätigkeitsbereich, welcher insbesonders die Pflege von
 von Patienten vorsieht, sowie einen mitverantwortlich en
 Tätigkeitsbereich, welcher die Durchführung von einfachen
 ärztlichen Tätigkeiten (Blutentnahme, Verabreichung von
Infusionen und Medikamenten, \ldots) vorsieht. 
Für spezielle Tätigkeitsbereiche gibt es entsprechende
Sonderausbildungen (OP-Bereich, Anästhesie, Intensivpflege,
\ldots)
\jNEUende

\section{[RSP]}
Überarbeitung SAMPLER:
\subsection{Atemwegsverlegung}
 \Layout{0}
 {\TxABCDE}
 {~
 \begin{ABCDEText}
   \item[\ABCDE{1}] Achte auf Essenreste, Spielsachen
   (Murmeln!) und andere potentielle Fremdkörper.
   %%%
   \item[\ABCDE{2}] Oft gestikuliert der Patient
   entsprechend. Bei fortgeschrittener Hypoxie kann der
   Patient {\RedFlag}\,zyanotisch sein.
   %%%
   \item[\ABCDE{3}] Infolge einer Hypoxie
   kann es zu einer {\RedFlag}\,\E{Bewusstseinstrübung} kommen.
   %%%
   \item[\ABCDE{4}] Das Leitsymptom ist Atemnot bzw.
   Fremdkörpergefühl, zusammen mit einer entsprechenden
   Anamnese.
   %%%
   \item[\ABCDE{A}] Der Atemweg ist teilweise oder ganz
   verlegt. Eventuell ist ein pfeifendes Atemgaeräusch
   (\I{Stridor}) zu hören. Bei einer hochgradigen
   oder kompletten Verlegung ist das Atemgeräusch
   abgeschwächt oder fehlt völlig. Der Fremdkörper ist oft
  nicht sichtbar, da er sich oft unterhalb des Rachenraums
   befindet. 
   %%%
   \item[\ABCDE{B}] Je nach Ausmaß der Verlegung kann die
  Sauerstoffsättigung erniedrigt sein. Der Patient \deleted{weisst} \added{hat}
  eine schnelle und erschwerte Atmung \deleted{auf} (\I{Tachypnoe}).
   \begin{itemize}
   \item \T[Atemwegsverlegung!milde]{Milde Verlegung}: Bei
   der milden Verlegung kann der Patient noch sprechen und effektiv husten. Er klagt
   meist über Atemnot und ein Fremdkörpergefühl, oft versucht er zu husten
   oder zu würgen. Oft ist auch ein pfeifendes Atemgeräusch
   während der Ein- und/oder Austamung zu hören
   (in-/exspiratorischer \E{Stridor}).
   \item \T[Atemwegsverlegung!schwere]{Schwere Verlegung}:
   Bei einer schweren Verlegung der Atemwege ist der Patient
  \deleted{\EE{unfähig zu reden oder effektiv zu husten}. Eventuell kann
  er nur durch Nic\-k\-en oder Gesten kommunizieren. In besonders
  schweren Fällen kann es zu einem \EE{Atemstopp} und nach
  kurzer Zeit zu einer \EE{Bewusstlosigkeit}, gefolgt von einem
  \EE{Kreislaufstillstand} kommen.}
  \added{{\RedFlag}\,\EE{unfähig zu reden oder effektiv zu husten}.
  Eventuell kann er nur durch Nic\-k\-en oder Gesten kommunizieren. In besonders
  schweren Fällen kann es zu einem
  {\RedFlag}\,\EE{Atemstopp} und nach kurzer Zeit zu einer
  {\RedFlag}\,\EE{Bewusstlosigkeit}, gefolgt von einem \EE{Kreislaufstillstand}
kommen.}
 \end{itemize}
   %%%
  \item[\ABCDE{C}] \deleted{Der Patient ist tachykard.}
  %%%
  \item[\ABCDE{D}] \deleted{Oft liegt eine Hypertonie aufgrund des}
  \item[\ABCDE{C}] \added{Der Patient ist tachykard. Oft liegt eine Hypertonie
aufgrund des Stress vor.}
  %%%
%   \item[\ABCDE{D}]
 %   %%%
 %   \item[\ABCDE{E}] 
 %   %%%
 %   \item[\ABCDE{\ldots}] 
 %   %%%
   \item[\ABCDE{=}] Bei einer schweren Atemwegsverlegung ist
   der Patient vital bedroht. Bei der milden Verlegung ist
   die Einschätzung von den Befunden abhängig.
 \end{ABCDEText}
 }
 {
 \begin{ABCDESlide}
  \item[\ABCDE{1}] \added{Essensreste, Spielsachen, etc.}
  \item[\ABCDE{2}] \added{Gestik, evtl. {\RedFlag}\,Zyanose}
  \item[\ABCDE{3}] \added{Evtl. {\RedFlag}\,\E{Bewusstseinstrübung}}
  \item[\ABCDE{4}] %
  \begin{itemize}
    \item \deleted{Atemnot}
    \item \deleted{Fremdkörpergefühl}
  \end{itemize}
  \item[\ABCDE{4}] \added{Atemnot, Fremdkörpergefühl}
   %%%
   \item[\ABCDE{A}] %
   \begin{itemize}
     \item Fremdkörper oft nicht sichtbar
     \item Evtl. pfeifendes Atemgeräusch (Stridor)
   \end{itemize}
   %%%
   \item[\ABCDE{B}] %
   \begin{itemize}
     \item Atemnot
    \item \added{Milde Verlegung: \EE{Kann Sprechen \& Husten}}
     \item Schwere Verlegung:
     \begin{itemize}
      \item \deleted{Unfähigkeit zu Sprechen oder effektiv zu husten}
      \item \deleted{Evtl. Bewusstlosigkeit → Kreislaufstillstand}
    \end{itemize}
    \item \deleted{Milde Verlegung:}
    \begin{itemize}
      \item \deleted{\EE{Kann Sprechen \& Husten}}
      \item \added{{\RedFlag}\,Unfähigkeit zu Sprechen oder
      effektiv zu husten}
      \item \added{Evtl. {\RedFlag}\,Bewusstlosigkeit,
      Hypoxie bis Kreislaufstillstand möglich}
     \end{itemize}
   \end{itemize}
  \item[\ABCDE{C}] \added{HF $\uparrow$, RR $\uparrow$}
  \item[\ABCDE{=}] \added{Bei schwerer Verlegung vital bedroht,
  sonst je nach Befund}
 \end{ABCDESlide}
 }
 {}
 {}
 {}

\subsection{Asthma}
\jNEU
\Layout{0}{Status Asthmaticus}{%
\index{Status!asthmaticus}
\label{topic:status-asthmaticus}
Der
\I[Status!asthmaticus]{Status asthamticus} ist eine gefürchtete Komplikation des
Asthmaanfalles bei der die Atemnot nicht wieder vergeht. Es
besteht die Gefahr, dass der Patient nach einiger Zeit keine
Kraft mehr zum Atmen hat (Erschöpfung). Der Status
asthmaticus kann tödlich enden.
}{
\begin{itemize}
  \item Lebensgefährliche Komplikation:
  \item Anfall hört nicht mehr auf
\end{itemize}
}{}{}{}
.
\jNEUende

 \Layout{0}
 {\TxABCDE}
 {~
 \begin{ABCDEText}
   \item[\ABCDE{1}] Kalte Temperaturen oder Allergene
   (Gräser, Pollen, \ldots) begünstigen einen Asthma-Anfall.
   %%%
   \item[\ABCDE{2}] Der Patient nimmt meist automatisch und unbewusst eine
   Position ein, die seine Atemhilfsmuskulatur unterstützt: Er
   sitzt aufrecht und stützt sich oft nach vorne oder hinten
   ab.
   %%%
   \item[\ABCDE{3}] Im schweren Anfall kann es zu
   {\RedFlag}\,\E{Bewusstseinsstörungen} kommen.
   %%%
   \item[\ABCDE{4}] Das Leitsymptom ist die Atemnot.
   %   %%%
   %   \item[\ABCDE{A}]
   %%%
   \item[\ABCDE{B}] Die Atemfrequenz ist erhöht
   (\E{Tachypnoe}), die {\RedFlag}\,\E{Sauerstoffsättigung}
   kann erniedrigt sein. Bei einer manifesten Hypoxie kann eine
   {\RedFlag}\,\E{Zyanose} beobachtbar sein.
   
   Oft ist ein pfeifendes,
   brummendes oder giemendes Atemgeräusch, besonders während
   der Ausatmung (\E{Exspiration}, \E{exspiratorischer
   Stridor}) wahrnehmbar. Die Exspiration ist ausserdem
   verlängert. Trockener Husten ist häufig.
   
  \deleted{\Z{Unterschätze nie die Schwere eines Asthmaanfalles!
  Gefährlich sind besonders die \protect\enquote*{ruhigen},
  wenig klagenden Patienten, diese haben oft ein anderes
  Luftnot-Empfinden. Diese Patienten sind Hochrisiko-Patienten
  und haben auch mit der \ce{O2}-Gabe Probleme! \ce{O2}-Gabe nur
  bis max. 95\PC* Sauerstofsättigung!
  \citep{Olschewski2009}\A{\footnote{Vortrag Univ-Prof.Dr. Horst Olschewski, 5. AGN-Kongress 2009; Notfall Rettungsmed (2009) 12:322}}}
}   
   Im \EE{schweren Anfall} ist oft \E{kein}, oder nur mehr ein
   \E{leises Atemgeräusch} hörbar, da kaum mehr Luft bewegt
   wird \Z{(\I{silent chest})}. Die
   {\RedFlag}\,Sauerstoffsättigung ist stark vermindert und der
   Patient ist {\RedFlag}\,\E{zyanotisch}. Im Endstadium kann
  die {\RedFlag}\,\E{Atemarbeit erschöpft} sein. 
  
  \jNEU
  Bei der Einschätzung der Atmung ist es wichtig auf die
  \EE{Zeichen einer erschwerten Atemarbeit} zu
  achten:
  \begin{itemize}
    \item Einsatz der Atemhilfsmuskulatur
    \item Einziehungen an den Rippen
    \item Erschöpfungszeichen
  \end{itemize}
  Es besteht die Gefahr der {\RedFlag}\,\EE{Erschöpfung der
  Atemarbeit}.
  \jNEUende
   %%%
   \item[\ABCDE{C}] Die Herzfrequenz ist meist erhöht
   (\E{Tachykardie}).
  \item[\ABCDE{=}] \added{Bei schwerer Atemnot besteht 
  eine vitale Bedrohung. Bei Erschöpfung der Atemarbeit
  besteht akute Lebensgefahr!}
 \end{ABCDEText}
 }
 {
 \begin{ABCDESlide}
  \item[\ABCDE{1}] \added{Kälte, Allergene}
   %%%
   \item[\ABCDE{2}] Einsatz der Atemhilfsmuskulatur
   %%%
   \item[\ABCDE{3}] Schwerer Anfall:
   {\RedFlag}\,Bewusstseinsstörungen
   %%%
   \item[\ABCDE{4}] Akute Atemnot
 %   %%%
 %   \item[\ABCDE{A}] 
   %%%
   \item[\ABCDE{B}] Akute Atemnot mit trockenem
   Husten, Verlängerte Ausatmungsphase, AF $\uparrow$,
   Exspiratorisches Atemgeräusch (Stridor)
   
   Schwerer Anfall: Kein oder leises Atemgeräusch,
   {\RedFlag}\,Sp\ce{O2} $\downarrow\downarrow$

\jNEU
  {Zeichen einer erschwerten Atemarbeit}:
  \begin{itemize}
    \item Einsatz der Atemhilfsmuskulatur
    \item Einziehungen an den Rippen
    \item Erschöpfungszeichen
  \end{itemize}
  .
  \jNEUende
  
   Endstadium: {\RedFlag}\,Atemarbeit erschöpft
   %%%
   \item[\ABCDE{C}] HF $\uparrow$
   
   Schwerer Anfall: HF $\downarrow$
   %%%
   \item[\ABCDE{D}] 
   Schwerer Anfall: RR $\downarrow$
   %%%
  \item[\ABCDE{=}] \added{Schwere Atemnot: Vital bedroht.
  Erschöpfte Atemarbeit: Akute Lebensgefahr}
 \end{ABCDESlide}
 }{}{}{}
 .
 
 \jALT
 .
 \subparagraph{\Ca{Schwerer Anfall}}
 \ldots
 \jALTende
 
 \jNEU
  \Layout{0}
 {{\TxSAMPLER}}
{%
 \begin{SAMPLERText}
  \item[\SAMPLER{A}] Patienten mit allergischem Asthma haben
  häufig Allergien gegen Gräser, Pollen, Tierhaare usw.
  \item[\SAMPLER{M}] Oft werden Sprays bzw. Inhalatoren zur
  Dauertherpaie und für den Anfall eingesetzt. Manche
  Medikamente, insbesonders Schmerzmedikamente, können
  Anfälle auslösen.%
  \ZFootnote{Besonders nichtsteroidale Antirheumatika (ASS,
  Aspirin{\texttrademark}, etc.) und Betablocker
  (Antiarrhythmika, Blutdruckmittel) können Asthma-Anfälle
  auslösen.} 
  Nahm der Patient in letzter Zeit (neue)
  Medikamente?
  \item[\SAMPLER{P}] Meist ist Asthma vorbekannt. 
  \item[\SAMPLER{L}] Letzte Spray-Einnahme erfragen!
  \item[\SAMPLER{E}] Häufige Auslösefaktoren sind Stress,
  Anstrengung und Allergenexposition.
%   \item[\SAMPLER{R}]
 \end{SAMPLERText}
 }
 {
 \begin{SAMPLERSlide}
%   \item[\SAMPLER{S}]
  \item[\SAMPLER{A}] Gräser, Pollen, Tierhaare, \ldots
  \item[\SAMPLER{M}] Sprays, Inhalatoren. (Neue) Medikamente in letzter Zeit?
  \item[\SAMPLER{P}] Asthma oft vorbekannt
  \item[\SAMPLER{L}] Spray-Einnahme?
  \item[\SAMPLER{E}] Stress, Anstrengung, Allergene
%   \item[\SAMPLER{R}]
 \end{SAMPLERSlide}
 }
 {}
 {}
 {}
 .
 \jNEUende
 
 \jALT
 \Layout{0}{Status Asthmaticus}{%
\index{Status!asthmaticus}
\label{topic:status-asthmaticus}
Der
\protect\enquote{Status} ist eine gefürchtete Komplikation des
Asthmaanfalles bei der die Atemnot nicht wieder vergeht. Es
besteht die Gefahr, dass der Patient nach einiger Zeit keine
Kraft mehr zum Atmen hat (Erschöpfung). Der Status
asthmaticus kann tödlich enden.
}{
\begin{itemize}
  \item Lebensgefährliche Komplikation:
  \item Anfall hört nicht mehr auf
\end{itemize}
}{}{}{}
 .
 \jALTende
 
\subsection{COPD}
\jALT
\subsubsection{Chronische Bronchitis und 
Chronische Obstruktive Lungenerkrankung
(COPD)}

 \DefinitionInPara{Die \T{chronische Bronchitis} ist
 eine chronische entzündliche Schleimhautschädigung  der unteren Atemwege.}
 \DefinitionInPara{Die \T{COPD}
 (\T[Lungenerkraknkung!chronisch-obstruktive]{Chronische Obstruktive
 Lungenerkrankung}\ZFootnote{\Lang{engl.} Chronic
 obstructive pulmonary disease\index{Chronic
 obstructive pulmonary disease|Z})} ist eine chronische entzündliche
 Schleimhautschädigung, welche eine zunehmende obstruktive 
 Atemwegseinschränkung aufweist.}
 Am Anfang steht die chronische Bronchitis, welche durch
 Husten mit schleimigem Auswurf gekennzeichnet ist
({\I{Raucherhusten}}). Es kommt dabei zu einer gesteigerten
 Entzündungsantwort auf eingeatmete Stoffe (Zigarettenrauch,
 Umweltschadstoffe, \ldots). Wenn man in der
Lungenfunktionsuntersuchung eine Atemwegseinschränkung
 nachweisen kann, spricht man von der chronisch-obstruktiven
 Lungenerkrankung (\Abk{COPD}). 
 Diese geht mit bleibenden Veränderungen der unteren Atemwege einher, 
 diese sind in der Grafik  \prettyref{fig:copd-veraenderungen} dargestellt.
 
% \begin{figure}[H]
% 
% \includegraphics[width=\linewidth]{./images/hirtler-aass/Alveolen-Bronchus-COPD-edited2.png}
% \Captionof{figure}{\label{fig:copd-veraenderungen}Veränderung der
% Atemwege bei der COPD. \TabSub{Links:} Schema eines gesunden
% Bronchus und einer gesunden Alveole. \TabSub{Oben rechts:} Bei der COPD sind
% die kleinen Luftwege verschleimt und verengt. \TabSub{Unten rechts:} Die
% Lungenbläschen (Alveolen) sind überbläht, weil die Luft nur erschwert wieder entweichen kann.
% \SourcePicture{Hirtler.}}
% \end{figure}



\Layout{0}{\TxSymptome}{%
Die \Abk{COPD} ist durch eine voranschreitende
Verschlechterung der Atemleistung gekennzeichnet, erstes
Anzeichen ist ein Husten mit Auswurf, der klassische
\protect\enquote{Raucherhusten}.
Je nach Schweregrad kommt es zu Zeichen der
Atemwegsverlegung (Obstruktion) und Ateminsuffizienz: Die
\E{Ausatmung} (Exspiration) ist erschwert, der Patient klagt
über \E{Belastungsdyspnoe}, ein Engegefühl, \E{Husten} und
hat Tachypnoe und evtl. Zyanose. Man kann oft ein
\E{brummendes, pfeifendes oder giemendes Atemgeräusch}
hören.

 \Z{Durch die erschwerte Ausatmung kommt es zu einer
 chronischen Überblähung der Lungenbläschen und zu einer \protect\enquote{Faß-förmigen}
 Verformung des Brustkorbes. Im Endstadium zeigen sich
 Zeichen einer Rechtsherzinsuffizienz
 (\prettyref{topic:rechtsherzinsuffizienz}) aufgrund einer
 Störung im Lungenkreislauf.}
 
Die \EE{Sauerstoffsättigung} bei einem COPD-Patienten ist
auch ohne akute Symptome meist \E{deutlich vermindert}
($\sim$\,{90}\PC*). Der sonst übliche Normalwert
gilt hier nicht! Bei fortgeschrittener Ateminsuffizienz
bekommt der Patient oft eine \E{Heimsauerstofftherapie} verordnet.
 }{
 \begin{itemize}
   \item Husten mit Auswurf, anfänglich
   \protect\enquote*{Raucherhusten}
   \item Voranschreitende, sich verschlechternde \E{Ateminsuffizienz}:
   \begin{compactitem}
     \item Belastungsdyspnoe (bei schwererer \Abk{COPD}
     auch in Ruhe)
     \item Ausatmung (Exspiration) erschwert
     \item Brummendes, pfeifendes od. giemendes Ausatemgeräusch
     \item Evtl. chronische Zyanose
     \item Evtl. Heimsauerstoff
     \item Sp\ce{O2} chronisch niedrig
   \end{compactitem}
\end{itemize}
}{}{}{}



\Layout{0}{Exazerbationen}{%
Der \Abk{COPD}-Patient erleidet öfters plötzliche
Verschlechterungen, sog. \E{Exazerbationen}. Diese sind
meist durch Infekte ausgelöst und durch vermehrte Atemnot,
Husten und Auswurf gekennzeichnet.
}{
 \begin{itemize}
   \item Plötzliche Verschlechterungen
   \item Meist Infekt-bedingt
   \item Vermehrte Atemnot,
   Husten und Auswurf
 \end{itemize}
 }{}{}{} 
 
%  \begin{figure}[h]
%  \begin{center}
%  % \includegraphics{./images/copd-stufen.png}
%  \includegraphics[width=\linewidth]{./images/hirtler-aass/COPD-Stufen_edited.pdf}
%  
%  \Captionof{figure}{Eine COPD
%  entsteht nicht plötzlich: Ein COPD{\textquotesingle}ler hat eine
%  \protect\enquote{Karriere} hinter sich. \SourcePicture{Hirtler}}
%  \end{center}
%  \end{figure}
 
 
\Layout{0}{Probleme mit Sauerstoff bei
COPD-Patienten}
{
Normalerweise wird der Atemantrieb über die
\ce{CO2}-Konzentration im Blut gesteuert. Die
\ce{CO2}-Konzentration ist bei \Abk{COPD}-Patienten jedoch chronisch
hoch, der Körper gewöhnt sich an diesen Zustand. Er
regelt dann den Atemantrieb direkt über den \ce{O2}-Gehalt im Blut.
Wird der \ce{O2}-Gehalt plötzlich künstlich erhöht, fehlt
der \index{COPD!Atemantrieb}Atemantrieb und der Patient kann
\EE{Bewusstseinsstörungen} (\protect\enquote{\ce{CO2}-Narkose}) und einen
\EE{Atemstillstand} erleiden.

}{
\Ca{Achtung!}

\SymbolText
}{}{}{}
.
\jALTende

\jNEU
\subsubsection{Chronische
Bronchitis und COPD}

\Layout{0}{\TxBeschreibung: COPD}{%
 \DefinitionInPara{Die \T{chronische Bronchitis} ist
 eine chronische entzündliche Schleimhautschädigung  der unteren Atemwege.}
 \DefinitionInPara{Die \T{COPD}
 (\T[Lungenerkraknkung!chronisch-obstruktive]{Chronische Obstruktive
 Lungenerkrankung}\ZFootnote{\Lang{engl.} Chronic
 obstructive pulmonary disease\index{Chronic
 obstructive pulmonary disease|Z})} ist eine chronische entzündliche
 Schleimhautschädigung, welche eine zunehmende obstruktive 
 Atemwegseinschränkung aufweist.}
Die \Abk{COPD} ist durch eine voranschreitende
Verschlechterung der Atemleistung gekennzeichnet.
 Am Anfang steht die chronische Bronchitis, welche durch
 Husten mit schleimigem Auswurf gekennzeichnet ist
(\enquote*{\I{Raucherhusten}}). Es kommt dabei zu einer gesteigerten
 Entzündungsantwort auf eingeatmete Stoffe (Zigarettenrauch,
 Umweltschadstoffe, \ldots). Wenn man in der
Lungenfunktionsuntersuchung eine Atmungsseinschränkung
 nachweisen kann, spricht man von der chronisch-obstruktiven
 Lungenerkrankung (\Abk{COPD}). 
 Diese geht mit bleibenden Veränderungen der unteren Atemwege einher, 
 diese sind in der Grafik  \prettyref{fig:copd-veraenderungen} dargestellt.
 
 \Z{Durch die erschwerte Ausatmung kommt es zu einer
 chronischen Überblähung der Lungenbläschen und zu einer \protect\enquote{Faß-förmigen}
 Verformung des Brustkorbes. Im Endstadium zeigen sich
 Zeichen einer Rechtsherzinsuffizienz
 (\prettyref{topic:rechtsherzinsuffizienz}) aufgrund einer
 Störung im Lungenkreislauf.}
 
\subparagraph{Exazerbationen}
Kommt zu der ohnehin schweren Grunderkrankung noch ein
erschwerender Faktor hinzu, z.\,B. eine Infektion der
Atemwege, kann es zu einer plötzlichen Verschlechterung
kommen, zur \E{Exazerbation}. Diese ist meist durch
vermehrte Atemnot, Husten und Auswurf gekennzeichnet.
 }{
 \begin{itemize}
   \item Husten mit Auswurf, anfänglich
   \protect\enquote*{Raucherhusten}
   \item Voranschreitende, sich verschlechternde \E{Ateminsuffizienz}:
   \begin{compactitem}
     \item Belastungsdyspnoe (bei schwererer \Abk{COPD}
     auch in Ruhe)
     \item Ausatmung (Exspiration) erschwert
     \item Brummendes, pfeifendes od. giemendes Ausatemgeräusch
     \item Evtl. chronische Zyanose
     \item Evtl. Heimsauerstoff
     \item Sp\ce{O2} chronisch niedrig
   \end{compactitem}
  \item Exazerbation:
 \begin{itemize}
   \item Plötzliche Verschlechterungen
   \item Meist Infekt-bedingt
   \item Vermehrte Atemnot,
   Husten und Auswurf
 \end{itemize}
\end{itemize}
 }{}{}{}
 
\Layout{0}{Probleme mit Sauerstoff bei
COPD-Patienten}
{
Normalerweise wird der Atemantrieb über die
\ce{CO2}-Konzentration im Blut gesteuert. Die
\ce{CO2}-Konzentration ist bei \Abk{COPD}-Patienten jedoch chronisch
hoch, der Körper gewöhnt sich an diesen Zustand. Er
regelt dann den Atemantrieb direkt über den \ce{O2}-Gehalt im Blut.
Wird der \ce{O2}-Gehalt plötzlich künstlich erhöht, fehlt
der \index{COPD!Atemantrieb}Atemantrieb und der Patient kann
\EE{Bewusstseinsstörungen} (\I{\ce{CO2}-Narkose}) und
im Extremfall einen \EE{Atemstillstand} erleiden.

}{
\Ca{Achtung!\\
Bewusstseinsstörungen bei übermäßiger
Sauerstoffgabe möglich!}

\SymbolText
}{}{}{}
 
 
 
% \begin{figure}[H]
%  
% \includegraphics[width=\linewidth]{./images/hirtler-aass/Alveolen-Bronchus-COPD-edited2.png}
% \Captionof{figure}{\label{fig:copd-veraenderungen}Veränderung der
% Atemwege bei der COPD. \TabSub{Links:} Schema eines gesunden
% Bronchus und einer gesunden Alveole. \TabSub{Oben rechts:} Bei der COPD sind
% die kleinen Luftwege verschleimt und verengt. \TabSub{Unten rechts:} Die
% Lungenbläschen (Alveolen) sind überbläht, weil die Luft nur erschwert wieder entweichen kann.
% \SourcePicture{Hirtler.}}
% \end{figure}

\Layout{0}
{\TxABCDE: COPD}
{~
\begin{ABCDEText}
  \item[\ABCDE{1}] Evtl. kalte Umgebung (kalte Luft führt
  zur Verengung der Bronchien)
  \item[\ABCDE{2}] Einsatz der Atemhilfsmuskulatur, Mühe
  beim Atmen, evtl. Heimsauerstoff, oft wirkt der Patient
  ängstlich.
  \item[\ABCDE{3}] Bei schweren Anfällen kann es aufgrund
  der Hypoxie zu {\RedFlag}\,\EE{Bewusstseinsstörungen}
  kommen.
  \item[\ABCDE{4}] Atemnot
  \item[\ABCDE{A}] Siehe \ABCDE{B}
  \item[\ABCDE{B}] Je nach Schweregrad kommt es zu Zeichen
  einer Atemwegsverlegung der \E{unteren}
  Atemwege  (Obstruktion, durch Verengung der Bronchien und
  Schleimproduktion) und Ateminsuffizienz:
  Die \E{Ausatmung} (Exspiration) ist erschwert, der Patient klagt
  über \E{Belastungsdyspnoe}, ein Engegefühl, \E{Husten} und
  hat Tachypnoe und evtl. {\RedFlag}\,\EE{Zyanose}. Man kann
  oft ein
  \E{brummendes, pfeifendes oder giemendes Atemgeräusch}
  hören.
  
  Die \EE{Sauerstoffsättigung} bei einem COPD-Patienten ist
  auch ohne akute Symptome meist \E{deutlich vermindert}
  ($\sim$\,{90}\PC*). Der sonst übliche Normalwert
  gilt hier nicht! Bei fortgeschrittener Ateminsuffizienz
  bekommt der Patient oft eine \E{Heimsauerstofftherapie}
  verordnet. Bei übermäßiger Sauerstoffgabe kann es zu
  {\RedFlag}\,\EE{Bewusstseinsstörungen kommen}
  (\I{\ce{CO2}-Narkose})!
  \item[\ABCDE{C}] Evtl. Tachykardie und hyperton
%   \item[\ABCDE{D}]
%   \item[\ABCDE{E}]
  \item[\ABCDE{=}] Bei {\RedFlag}\,schwerer Atemnot
  oder {\RedFlag}\,Bewusstseinsstörungen vitale Bedrohung.
  \item[\ABCDE{\ldots}] Als Nebenbefund können Infektzeichen
  vorliegen (erhöhte Körpertemperatur).
\end{ABCDEText}
}
{
\begin{ABCDESlide}
  \item[\ABCDE{1}] Kalte Umgebung
  \item[\ABCDE{2}] Atemhilfsmuskulatur, Anstrengung beim
  Atmen, ängstlich
  \item[\ABCDE{3}] Evtl. {\RedFlag}\,Bewusstseinsstörungen
  \item[\ABCDE{4}] Atemnot
  \item[\ABCDE{A}] Siehe \ABCDE{B}
  \item[\ABCDE{B}] Atemnot! Ausatmung erschwert, Husten
  Tachypnoe, {\RedFlag}\,Zyanose;
  
  Brummendes, pfeifendes Atemgeräusch
  
  Sp\ce{O2} chronisch niedrig ($\sim$\,{90}\PC*): Normalwert
  gilt nicht!
  \item[\ABCDE{C}] HF $\uparrow$, RR $\uparrow$
%   \item[\ABCDE{D}] 
%   \item[\ABCDE{E}] 
  \item[\ABCDE{=}] Vitale Bedrohung bei schwerer Atemnot und 
  Bewusstseinsstörungen
  \item[\ABCDE{\ldots}] Evtl. erhöhte Körpertemperatur
\end{ABCDESlide}
}
{}
{}
{}

\Layout{0}
{{\TxSAMPLER}: COPD}
{
\begin{SAMPLERText}
%   \item[\SAMPLER{S}] 
%   \item[\SAMPLER{A}]
  \item[\SAMPLER{M}] Inhalatoren bzw. Sprays zur Dauer- und
  Akuttherapie, Heimsauerstoff. \Z{Häufig Kortisonpräparate;
  bei vorbekanntem Infekt wurden oft schon Antibiotika
  verschrieben.}
  \item[\SAMPLER{P}] Eine COPD ist normalerweise vorbekannt.
  \item[\SAMPLER{L}] Letzte Spray-Einnahme?
%   \item[\SAMPLER{E}] 
%   \item[\SAMPLER{R}]
\end{SAMPLERText}
}
{
\begin{SAMPLERSlide}
%   \item[\SAMPLER{S}]
%   \item[\SAMPLER{A}]
  \item[\SAMPLER{M}] Inhalatoren, evtl. Heimsuaerstoff.
  \item[\SAMPLER{P}] COPD vorbekannt
  \item[\SAMPLER{L}] Letzte Spray-Einnahme?
%   \item[\SAMPLER{E}] 
%   \item[\SAMPLER{R}]
\end{SAMPLERSlide}
}
{}
{}
{}
 .
 
\jNEUende

\subsection{Lungenembolie}
\jALT
 
 \Layout{0}
{\TxABCDE}
 {~
 \begin{ABCDEText}
  \item[\ABCDE{Szeneüberblick}]
  \item[\ABCDE{Eindruck}] 
  \item[\ABCDE{Bewusstsein}]
  \item[\ABCDE{Hauptbeschwerde}] Das Leitsymptom ist die \EE{Atemnot}, eventuell in Kombination mit (atemabhängigen) Thorax- oder Rückenschmerzen. 
  \item[\ABCDE{A}] 
   \item[\ABCDE{B}] Meist atmet der Patient schnell (Tachypnoe) und klagt über
einen \EE{atemabhängigen Thoraxschmerz}.
  \item[\ABCDE{C}]  Daneben kommt es zu
einer \E{Tachykardie} und Angstgefühlen., oft findet man deutlich hervortretende, \EE{gestaute Venen}
am Hals aufgrund des Blutrückstaus. Evtl. kann man auch einen
\E{arrhythmischen Puls} aufgrund eines Vorhofflimmerns, ertasten.
  \item[\ABCDE{D}] 
  \item[\ABCDE{E}] 
  \item[\ABCDE{\ldots}] Je nach Ursache der Embolie und Herkunft des Embolus kann
man Zeichen anderer Krankheitsbilder wahrnehmen: Bei einer
\EE{tiefen Beinvenenthrombose} ist ein Bein geschwollen,
überwärmt, rot/rot-bläulich verfärbt und schmerzhaft
(\prettyref{topic:thrombose}). 

Eine Lungenembolie kann man leicht mit einem Herzinfarkt
oder Angina pectoris verwechseln. Dies hat aber bei der
Erstversorgung kaum Konsequenz, die Erstversorgen ist gleich.
  \item[\ABCDE{=}] 
 \end{ABCDEText}
 }
 {
 \begin{ABCDESlide}
  \item[\ABCDE{Sz}] 
  \item[\ABCDE{Ei}] 
  \item[\ABCDE{Be}] 
  \item[\ABCDE{Hb}] Dyspnoe, (atemabhängige) Thorax- od. Rückenschmerzen, Angstgefühl
  \item[\ABCDE{A}] 
  \item[\ABCDE{B}] Atemnot, Tachypnoe
  \item[\ABCDE{C}] Tachkardie, evtl. gestaute Halsvenen (Blutrückstau) , evtl. arrhythmischer Puls bei Vorhofflimmern
  \item[\ABCDE{D}] 
  \item[\ABCDE{E}] 
  Evtl. Bein: einseitig geschwollen,
  überwärmt, rot/rot-bläulich verfärbt, schmerzhaft (Beinvenenthombose)
  \item Evtl. arrhythmischer Puls (Vorhofflimmern)
  \item[\ABCDE{=}] 
 \end{ABCDESlide}
 }
 {}
 {}
 {}
 
 \Layout{0}
{{\TxSAMPLER}}
 {
 \begin{SAMPLERText}
  \item[\SAMPLER{S}] 
  \item[\SAMPLER{A}]
  \item[\SAMPLER{M}] 
   \item[\SAMPLER{P}]
Oft ist bei den Patienten auch \EE{Vorhofflimmern} als
chronische Erkrankung bekannt.zu leiden.
   \item[\SAMPLER{L}]
  \item[\SAMPLER{E}] 
   \item[\SAMPLER{R}] \EE{Immobilisation}, z.\,B. durch Bettlägrigkeit,
Liegegipsverbände oder lange Reisen, innerhalb der letzten 4
Wochen, erhöht das Risiko einer Thrombose und damit auch das
einer Lungenembolie. \Z{Weiters tritt die Lungenembolie gehäuft in der Schwangerschaft
auf. Rauchen sowie die Einnahme der \protect\enquote{Pille} erhöhen
ebenso das Risiko.}
 \end{SAMPLERText}
 }
 {
 \begin{SAMPLERSlide}
  \item[\SAMPLER{S}]
  \item[\SAMPLER{A}]
  \item[\SAMPLER{M}] 
  \item[\SAMPLER{P}] 
  \item[\SAMPLER{L}]
  \item[\SAMPLER{E}] 
  \item[\SAMPLER{R}]Immobilisation, Vorhofflimmern, \Z{Schwangerschaft}, \Z{Rauchen}, \Z{Verhütungspille}
 \end{SAMPLERSlide}
 }
 {}
 {}
 {}
. 

\jALTende

\jNEU
 
 \Layout{0}
{\TxABCDE: Lungenembolie}
 {~
 \begin{ABCDEText}
  \item[\ABCDE{1}] Hinweise auf Bettlägrigkeit
  (Pflegebett, \ldots) ? Hinweise auf eine Fernreise?
  \item[\ABCDE{2}] Hinweise auf Bettlägrigkeit (Alter,
  Gipsverband, \ldots)
%   \item[\ABCDE{3}]
  \item[\ABCDE{4}] Das Leitsymptom ist die
  {\RedFlag}\,\EE{Atemnot}, eventuell in Kombination mit
  (atemabhängigen) {\RedFlag}\,Thorax- oder Rückenschmerzen. 
  
  Eine Unterscheidung zwischen einer Lungenembolie und
  einem Akuten Koronarsyndrom ist oft nicht möglich. 
%   \item[\ABCDE{A}]
   \item[\ABCDE{B}] Meist atmet der Patient schnell (Tachypnoe) und klagt über
  einen \EE{atemabhängigen {\RedFlag}\,Thoraxschmerz}.
  \item[\ABCDE{C}]  Es kommt zu
  einer \E{Tachykardie} und Angstgefühlen, oft findet man
  deutlich hervortretende, \EE{gestaute Venen}
  am Hals aufgrund des Blutrückstaus. Evtl. kann man auch einen
  \E{arrhythmischen Puls} aufgrund eines Vorhofflimmerns, ertasten.
  Der Blutdruck kann normal oder vermindert sein. Bei einer
  schweren Lungenembolie kann es zu einem
  {\RedFlag}\,\EE{kardiogenen Schock} kommen, es ist daher
  wichtig auf {\RedFlag}\,Schockzeichen zu achten!
%   \item[\ABCDE{D}]
%   \item[\ABCDE{E}]
  \item[\ABCDE{=}] Bei {\RedFlag}\,bestehender Atemnot oder
  {\RedFlag}\,Thoraxschmerzen besteht eine vitale Bedrohung.
  \item[\ABCDE{\ldots}] Je nach Ursache der Embolie und
  Herkunft des Embolus kann
  man Zeichen anderer Krankheitsbilder wahrnehmen: Bei einer
  \EE{tiefen Beinvenenthrombose} ist ein Bein geschwollen,
  überwärmt, rot/rot-bläulich verfärbt und schmerzhaft
  (\prettyref{topic:thrombose}).
  
 \end{ABCDEText}
 }
 {
 \begin{ABCDESlide}
  \item[\ABCDE{1}] Bettlägrigkeit?
  \item[\ABCDE{2}] Bettlägrigkeit?
%   \item[\ABCDE{3}] 
  \item[\ABCDE{4}] Dyspnoe, (atemabhängige)
  {\RedFlag}\,Thorax- od. Rückenschmerzen, Angstgefühl
  
  Unterscheidung zu ACS oft nicht möglich
%   \item[\ABCDE{A}] 
  \item[\ABCDE{B}] {\RedFlag}\,Atemnot, Tachypnoe
  \item[\ABCDE{C}] HF $\uparrow$, evtl. arrhythmischer Puls
  (Vorhofflimmern), RR \SymbolFindingNormal od.
  {\RedFlag}\,$\downarrow$
 
  Gestaute Halsvenen (Blutrückstau), evtl. arrhythmischer Puls bei Vorhofflimmern
  
  {\RedFlag}\,Schockzeichen?
% 
%   \item[\ABCDE{D}] 
%   \item[\ABCDE{E}] 
  \item[\ABCDE{=}] Bei bestehender Atemnot oder
  Thoraxschmerzen vitale Bedrohung.
  \item[\ABCDE{\ldots}] Venöser Verschluss?
 \end{ABCDESlide}
 }
 {}
 {}
 {}
 
 \Layout{0}
{{\TxSAMPLER}: Lungenembolie}
 {
 \begin{SAMPLERText}
  \item[\SAMPLER{M}] Bettlägrigen Patienten werden zur
  Thromboseprophylace oft gerinnungshemmende Medikamente
  verschreiben (\I{Thrombosespritzen}). Wurden diese
  genommen?
  \item[\SAMPLER{P}]
  Oft ist bei den Patienten \EE{Vorhofflimmern} als
  chronische Erkrankung bekannt. Patienten, welche bereits
  schon einmal eine Thrombose oder eine Embolie hatten
  haben ein erhöhtes Risiko neuerlich daran zu erkranken. Im
  Rahmen von Krebserkrankungen kann es gehäuft zu Thrombosen
  und Embolien kommen.
   \item[\SAMPLER{L}]
  \item[\SAMPLER{E}] Fernreise? Langes Sitzen/Liegen? Wenig
  getrunken?
   \item[\SAMPLER{R}] \EE{Immobilisation}, z.\,B. durch Bettlägrigkeit,
  Gipsverbände oder lange Reisen, innerhalb der letzten 4
  Wochen, erhöht das Risiko einer Thrombose und damit auch das
  einer Lungenembolie. \Z{Weiters tritt die Lungenembolie gehäuft in der Schwangerschaft
  auf. Rauchen sowie die Einnahme der \protect\enquote{Pille} erhöhen
  ebenso das Risiko.}
 \end{SAMPLERText}
 }
 {
 \begin{SAMPLERSlide}
  \item[\SAMPLER{M}] Thrombosespritzen
  \item[\SAMPLER{P}] Vorhofflimmern, Thrombose, Embolie,
  Krebserkrankungen
%   \item[\SAMPLER{L}]
  \item[\SAMPLER{E}] Fernreise? Langes Sitzen/Liegen? Wenig
  getrunken?
  \item[\SAMPLER{R}] Immobilisation, Vorhofflimmern,
  \Z{Schwangerschaft}, \Z{Rauchen}, \Z{Verhütungspille}
 \end{SAMPLERSlide}
 }
 {}
 {}
 {}
 .
\jNEUende


\subsection{Lungenöden}
\jALT

\Layout{0}{\TxSymptome}
{%
\index{Brodelnde Atemgeräusche}%
Das Leitsymptom ist das \EE{brodelnde Atemgeräusch}, welches
in schweren Fällen bereits ohne Hilfsmittel hörbar und
typisch für das Lungenödem ist. Es klingt, als würde jemand
mit einem Strohhalm in ein Glas Wasser hineinblasen. Im
Extremfall kann man es bereits \E{auf mehrere Meter
Entfernung} wahrnehmen.
Weitere häufige Symptome sind \EE{Atemnot} und \EE{Zyanose}.
Der Patient bekommt oft nur \E{aufrecht sitzend} genügend
Luft (Orthopnoe).
}{
\begin{compactitem}
  \item Brodelndes Atemgeräusch
  \item Atemnot
  \item Zyanose
  \item Patient bekommt nur im Sitzen Luft
\end{compactitem}
}{}{}{}

 \Layout{0}
{\TxABCDE}
 {~
 \begin{ABCDEText}
  \item[\ABCDE{Szeneüberblick}]
  \item[\ABCDE{Eindruck}]
  \item[\ABCDE{Bewusstsein}]
  \item[\ABCDE{Hauptbeschwerde}]
  \item[\ABCDE{A}]
  \item[\ABCDE{B}]
  \item[\ABCDE{C}]
  \item[\ABCDE{D}]
  \item[\ABCDE{E}]
  \item[\ABCDE{=}] Patienten mit einem deutlich hörbaren Lungenödem (v.\,a. bei
  kardialer Ursache) sind schwer und \EE{lebensbedrohlich
   krank}. Das Herz verbraucht oft schon die allerletzten Reserven. In
   schweren Fällen kann schon die Aufregung durch den Transport
   im Stiegenhaus zum Auto zu einem Herz-Kreislaufstillstand
   führen.
   
   Es kommt tatsächlich relativ häufig vor, dass diese
   Patienten \EE{im Stiegenhaus} oder \EE{im Fahrzeug}
   \EE{reanimationspflichtig} werden. \Ca{Solche Patienten müssen
   daher \EE{an Ort und Stelle} eine entsprechende (ärztliche)
   Therapie erhalten bevor sie transportiert werden können.}
   
   Bei leichten Formen des Lungenödems sind die
   Atemgeräusche nicht mit freiem Ohr, sondern
   nur mittels Stethoskop zu hören.  Die Patienten zeigen auch
   sonst keine Alarmsymptome, sondern nur Ödeme und eine mäßige
   Kurzatmigkeit. Diese Patienten sind i.\,d.\,R. nicht
   unmittelbar vital bedroht, ein bestmögliches Monitoring ist jedoch wichtig.
 \end{ABCDEText}
 }
 {
 \begin{ABCDESlide}
  \item[\ABCDE{Sz}]
  \item[\ABCDE{Ei}]
  \item[\ABCDE{Be}]
  \item[\ABCDE{Hb}]
  \item[\ABCDE{A}]
  \item[\ABCDE{B}]
  \item[\ABCDE{C}]
  \item[\ABCDE{D}]
  \item[\ABCDE{E}]
   \item[\ABCDE{=}]
   \begin{itemize}
    \item Hörbares Lungenödem: \Ca{Lebensgefahr!}
     \item Leichte Form: Keine Alarmsymptome
   \end{itemize}
 \end{ABCDESlide}
 }
 {}
 {}
 {}
 
 \Layout{0}
{{\TxSAMPLER}}
 {
 \begin{SAMPLERText}
  \item[\SAMPLER{S}] 
  \item[\SAMPLER{A}]
  \item[\SAMPLER{M}] 
  \item[\SAMPLER{P}]
  \item[\SAMPLER{L}]
  \item[\SAMPLER{E}] 
  \item[\SAMPLER{R}]
 \end{SAMPLERText}
 }
 {
 \begin{SAMPLERSlide}
  \item[\SAMPLER{S}]
  \item[\SAMPLER{A}]
  \item[\SAMPLER{M}] 
  \item[\SAMPLER{P}] 
  \item[\SAMPLER{L}]
  \item[\SAMPLER{E}] 
  \item[\SAMPLER{R}]
 \end{SAMPLERSlide}
 }
 {}
 {}
 {}
 
\Layout{0}{Einschätzung}{%


}{
}{}{}{}
.
\jALTende

\jNEU

 \Layout{0}
{\TxABCDE: Lungenödem}
 {~
 \begin{ABCDEText}
  \item[\ABCDE{1}] Toxisches Lungenödem: Besconders auf
  Umgebungssicherheit achten: Offene Lacke,
  Lösungsmittel, Dämpfe, Chemikalien, Warnschilder, \ldots.
  Im Zweifel Anforderung der Feuerwehr! 
  \item[\ABCDE{2}] Der Patient sitzt, oft sind die Pölster
  im Bett aufgetürmt.
%   \item[\ABCDE{3}] 
  \item[\ABCDE{4}] Atemnot
  \item[\ABCDE{A}] Es kann zu schaumigen Auswurf kommen.
  \item[\ABCDE{B}] \index{Brodelnde Atemgeräusche}%
  Das Leitsymptom ist das
  {\RedFlag}\,\EE{brodelnde Atemgeräusch}, welches in
  schweren Fällen bereits ohne Hilfsmittel hörbar und
  typisch für das Lungenödem ist. Es klingt, als würde jemand
  mit einem Strohhalm in ein Glas Wasser hineinblasen. Im
  Extremfall kann man es bereits \E{auf mehrere Meter
  Entfernung} wahrnehmen.
  Weitere häufige Symptome sind {\RedFlag}\,\EE{Atemnot} und
  {\RedFlag}\,\EE{Zyanose}.
  Der Patient bekommt oft nur \E{aufrecht sitzend} genügend
  Luft \Z{(Orthopnoe)}, im Liegen verschlechtert sich die
  Atemnot deutlich.
  \item[\ABCDE{C}] Die Herzfrequenz und der Blutdruck können
  in jede Richtung verändert sein (Herzrhythmusstörungen
  oder ein hypertensiver Notfall können ein kardiales
  Lungenödem auslösen!). Im Verlauf kann es zu einem
  {\RedFlag}\,\EE{kardiogenen Schock} kommen.
%   \item[\ABCDE{D}]
%   \item[\ABCDE{E}] 
  \item[\ABCDE{=}] Patienten mit einem {\RedFlag}\,deutlich
  hörbaren Lungenödem (v.\,a. bei kardialer Ursache) sind schwer und \EE{lebensbedrohlich
   krank}. Das Herz verbraucht oft schon die allerletzten Reserven. In
   schweren Fällen kann schon die Aufregung durch den Transport
   im Stiegenhaus zum Auto zu einem Herz-Kreislaufstillstand
   führen.
   
   Es kommt tatsächlich relativ häufig vor, dass diese
   Patienten \EE{im Stiegenhaus} oder \EE{im Fahrzeug}
   \EE{reanimationspflichtig} werden. \Ca{Solche Patienten müssen
   daher \EE{an Ort und Stelle} eine entsprechende (ärztliche)
   Therapie erhalten bevor sie transportiert werden können.}
   
   Bei leichten Formen des Lungenödems sind die
   Atemgeräusche nicht mit freiem Ohr, sondern
   nur mittels Stethoskop zu hören.  Die Patienten zeigen auch
   sonst keine Alarmsymptome, sondern nur Ödeme und eine mäßige
   Kurzatmigkeit. Diese Patienten sind i.\,d.\,R. nicht
   unmittelbar vital bedroht, ein bestmögliches Monitoring ist jedoch wichtig.
  \item[\ABCDE{\ldots}] Beidseitige Beinödeme sind typisch.
 \end{ABCDEText}
 }
 {
 \begin{ABCDESlide}
  \item[\ABCDE{1}] Umgebung sicher? Chemikalien?
  Warnschilder?
  \item[\ABCDE{2}] Sitzend, aufgetürmte Pölster
%   \item[\ABCDE{3}]
  \item[\ABCDE{4}] Atemnot
  \item[\ABCDE{A}] Evtl. schaumiger Auswurf
  \item[\ABCDE{B}] {\RedFlag}\,Brodelndes Atemgeräusch

  Atemnot (besonders im Liegen), evtl. {\RedFlag}\,Zyanose
  
  Atemnot im Liegen schlechter
  \item[\ABCDE{C}] HF $\uparrow\downarrow$, RR
  $\uparrow\downarrow$
  
  Evtl. kardiogener Schock
%   \item[\ABCDE{D}]
%   \item[\ABCDE{E}]
   \item[\ABCDE{=}]
   \begin{itemize}
    \item {\RedFlag}\,Hörbares Lungenödem:
    \Ca{Lebensgefahr!}
     \item Leichte Form: Keine Alarmsymptome
   \end{itemize}
  \item[\ABCDE{\ldots}] Beidseitige Beinödeme
 \end{ABCDESlide}
 }
 {}
 {}
 {}
 
 \Layout{0}
{{\TxSAMPLER}: Lungenödem}
 {
 \begin{SAMPLERText}
  \item[\SAMPLER{M}] Herzpatienten nehmen of zahlreiche
  Medikamente ein (Blutdruck, Entwässerung, Blutfette,
  \ldots).
  \item[\SAMPLER{P}] Häufig berichten die Patienten über
  Herz-Kreislauferkraknkungen wie z.\,B. Herzinsuffizienz,
  Z.\,n. Herzinfarkt, Bluthochdruck oder
  Herzklappenerkrankungen.
%   \item[\SAMPLER{L}]
  \item[\SAMPLER{E}] Expostion von chemischen Substanzen?
  Ertrinkungsunfall?
  
  Erschwerender Faktor beim vorbestehender Herzerkrankung
  (Infekt, \ldots)?
%   \item[\SAMPLER{R}]
 \end{SAMPLERText}
 }
 {
 \begin{SAMPLERSlide}
  \item[\SAMPLER{M}] Zahlreich
  \item[\SAMPLER{P}] Herz-Kreislauferkraknkungen: Herzinsuffizienz,
  Z.\,n. Herzinfarkt, Bluthochdruck,
  Herzklappenerkrankungen, \ldots
%   \item[\SAMPLER{L}]
  \item[\SAMPLER{E}] Exposition? Ertrinkungsunfall?
  Erschwerender Faktor beim vorbestehender Herzerkrankung
  (Infekt, \ldots)?
%   \item[\SAMPLER{R}]
 \end{SAMPLERSlide}
 }
 {}
 {}
 {}
 .

\jNEUende


\subsection{Hyperventilationstetanie}
\jALT
 
\Layout{0}{\TxSymptome}
{%
Die Patienten beklagen \EE{Atemnot}, obwohl die
\ce{O2}-Zufuhr nicht beeinträchtigt und das Blut maximal
gesättigt ist (Sp\ce{O2} von 100\PC*). Sie haben eine
\EE{tiefe und schnelle Atmung} (\E{Tachypnoe}).
Durch das Elektrolyt-Ungleichgewicht kommt es zuerst zu
\EE{Schwindelgefühlen}, etwas später zu einem \EE{Kribbeln} in den
Fingern. In weiterer Folge zeigen sich tonische \EE{Krämpfe} 
(\T{Hyperventilationstetanie}).
Klassisch ist die sog.
\T{Pfötchenstellung}: Die Arme sind
angezogen und Handhaltung erinnert an Pfoten eines Tieres.
Wird die Hyperventilation nicht rechtzeitig beendet, kann es
zur Eintrübung bis hin zur Bewusstlosigkeit kommen.
}{
\begin{itemize}
  \item Subjektiv: \emph{\protect\enquote{Atemnot}}, Sp\ce{O2} von 100\PC* 
  \item Tiefe, schnelle Atmung
  \item Schwindel, Bewusstseinsstörungen
  \item Kribbeln in den Fingern
  \item Tonische Krämpfe (\protect\enquote{Pfötchenstellung})
\end{itemize}
}{}{}{}
 
 

\Layout{0}{Differentialdiagnosen}
{%
Die psychisch bedingte Hyperventilation muß unbedingt von
der Hyperventilation aufgrund anderer Ursachen abgegrenzt
werden. Bei \EE{Ateminsuffizienz} oder \EE{durch
Erkrankungen verursachte Atemnot} kommt es natürlich auch
oft zu einer (kompensatorischen) Hyperventilation
(Herzinfarkt, Lungenembolie, Pneumothorax, Asthma-Anfall,
\ldots).

Bei einer stoffwechselbedingten \EE{Übersäuerung} (Azidose)
kommt es zu einer Hyperventilation um \ce{CO2} abzuatmen
(dadurch wird Säure im Körper abgebaut). Die
\E{Kussmaul'sche Atmung} beim \EE{diabetischen Koma} wäre
ein typisches Beispiel dafür
(\prettyref{topic:koma-diabetisches}).

}{
\begin{itemize}
  \item Atemnot mit körperlicher Ursache
  \item Ateminsuffizienz
  \item Stoffwechselbedingte Azidose (z.\,B. Diabetisches Koma)
\end{itemize}
}{}{}{}
.
\jALTende

\jNEU
  
% 1-E
\Layout{0}
{\TxABCDE: Hyperventilationssyndrom}
{~
\begin{ABCDEText}
  \item[\ABCDE{1}] Oft kommt es bei vornehmlich
  jugendlichen Besuchern von Pop-Konzerten u.\,ä. zu
  Erregungszuständen, welche zum Hyperventilationssysdrom
  führen können. Daneben kann jedoch jede Situation, welche
  zu erhöhtem Stress oder Erregung führt, kann ein
  entsprechender Auslöser sein (Lernstress, Beziehungsende,
  \ldots). Das Schaffen bzw. Aufsuchen einer ruhigen
  Umgebung ist wichtig!
  \item[\ABCDE{2}] Häufig kann man Zeichen psychischer
  Erregtheit wahrnehmen: Wildes gestikulieren, Weinkrämpfe
  usw.
  
  Klassisch ist die sog. \T{Pfötchenstellung}:
  Die Arme sind angezogen und Handhaltung erinnert an Pfoten eines Tieres.
  \item[\ABCDE{3}] Wird die Hyperventilation nicht
  rechtzeitig beendet, kann es zur Eintrübung bis hin zur Bewusstlosigkeit kommen.
  \item[\ABCDE{4}] Das Leitsymptom ist die \EE{Atemnot}.
  
  Durch
  das Elektrolyt-Ungleichgewicht kommt es zuerst zu
  \EE{Schwindelgefühlen}, etwas später zu einem
  \EE{Kribbeln} in den
  Fingern. In weiterer Folge zeigen sich tonische \EE{Krämpfe}
  (\T{Hyperventilationstetanie}).
  %   \item[\ABCDE{A}]
  \item[\ABCDE{B}]  Die Patienten beklagen \EE{Atemnot}, obwohl die
  \ce{O2}-Zufuhr nicht beeinträchtigt und das Blut maximal
  gesättigt ist (Sp\ce{O2} von 100\PC*). Sie haben eine
  \EE{tiefe und schnelle Atmung} (\E{Tachypnoe}).
  \item[\ABCDE{C}] Herzfrequenz und Blutdruck können in
  Folge der Erregung etwas erhöht sein.
  %   \item[\ABCDE{D}]
  \item[\ABCDE{E}] \EE{Ausschlißen von gefährlichen
  Differentialdiagnosen}: Die psychisch bedingte
  Hyperventilation muß unbedingt von der Hyperventilation 
  aufgrund anderer Ursachen abgegrenzt
  werden. Bei \EE{Ateminsuffizienz} oder \EE{durch
  Erkrankungen verursachte Atemnot} kommt es natürlich auch
  oft zu einer (kompensatorischen) Hyperventilation
  (Herzinfarkt, Lungenembolie, Pneumothorax, Asthma-Anfall,
  \ldots).
  
  Bei einer stoffwechselbedingten \EE{Übersäuerung} (Azidose)
  kommt es zu einer Hyperventilation um \ce{CO2} abzuatmen
  (dadurch wird Säure im Körper abgebaut). Die
  \E{Kussmaul'sche Atmung} beim \EE{diabetischen Koma} wäre
  ein typisches Beispiel dafür
  (\prettyref{topic:koma-diabetisches}).
  \item[\ABCDE{=}] Normalerweise keine vitale Bedrohung
  (Differentialdiagnosen bedenken!).
  %   \item[\ABCDE{\ldots}]
\end{ABCDEText}
}
{
\begin{ABCDESlide}
  \item[\ABCDE{1}] Konzerte, stressige bzw. erregende
  Situationen
  \item[\ABCDE{2}] Gestikulieren, Wienkrämpfe,
  Pfötchenstellung
  \item[\ABCDE{3}] Evtl. Bewusstseinsstörungen
  \item[\ABCDE{4}] Atemnot. Evtl. Kribbeln, Schwindel,
  Krämpfe
  %   \item[\ABCDE{A}]
  \item[\ABCDE{B}] Klagt über Atemnot, Tachypnoe, aber sonst
  keine physischen Anzeichen (keine Zyanose, Sp\ce{O2}
  \SymbolFindingNormal, \ldots)
  \item[\ABCDE{C}] HF und RR evtl. $\uparrow$
  %   \item[\ABCDE{D}]
  \item[\ABCDE{E}] Ausschlissen:
  \begin{itemize}
    \item Atemnot mit körperlicher Ursache
    \item Ateminsuffizienz
    \item Stoffwechselbedingte Azidose (z.\,B. Diabetisches Koma)
  \end{itemize}
  \item[\ABCDE{=}] Normal keine vitale Bedrohung
  (Differentialdiagnosen bedenken!).
  %   \item[\ABCDE{\ldots}]
\end{ABCDESlide}
}
{}
{}
{}
 
\Layout{0}
{{\TxSAMPLER}: Hyperventilationssyndrom}
{
\begin{SAMPLERText}
%   \item[\SAMPLER{S}] 
%   \item[\SAMPLER{A}] 
%   \item[\SAMPLER{M}] 
%   \item[\SAMPLER{P}] 
%   \item[\SAMPLER{L}] 
  \item[\SAMPLER{E}] Oft wird über stressige oder aufregende
  Erlebnisse berichtet (Trennung, Todesfälle, Lernstress,
  \ldots).
%   \item[\SAMPLER{R}] 
\end{SAMPLERText}
}
{
\begin{SAMPLERSlide}
%   \item[\SAMPLER{S}] 
%   \item[\SAMPLER{A}] 
%   \item[\SAMPLER{M}] 
%   \item[\SAMPLER{P}] 
%   \item[\SAMPLER{L}] 
  \item[\SAMPLER{E}] Stressige oder aufregende
  Erlebnisse
%   \item[\SAMPLER{R}] 
\end{SAMPLERSlide}
}
{}
{}
{}
.
\jNEUende

\section{[SAN]}
 \Layout{2}{Material und Personal}{
 \begin{itemize}
    \item \deleted{Mind. 2 Helfer}
    \item \added{Mindestens 2 Helfer, besser 3 Helfer}
     \item Schaufeltrage
     \item Gurte für Schaufeltrage
 \end{itemize}
 }{}{}{}{}

\jALT
 \Layout{2}{Vorgehen}{
 \begin{enumerate}
     \item \ldots
     \item \EE{Öffnen} der Trage: \ldots
     \item \EE{Aufschaufeln}:
     \begin{enumerate}
       \item 1. Methode: Fußseite offen, Patient wird von oben
    kommend wie auf eine Schere aufgeladen, die Trage unter ihm
    geschlossen. Kontrolle des sicheren Schlusses!
    
    Diese Methode ist besonders bei Becken- und
    Oberschenkelverletzungen sinnvoll und der zweiten
    Methode vorzuziehen da der Patient nicht am Becken
    bewegt wird.
    \item 2. Methode: beide Teile getrennt, Patient leicht an
    Schulter und Hüfte zur Seite gedreht und jeweiligen Tragenteil
    darunter geschoben. Hälften zuerst am Kopf-, dann am Fußende
    schließen. Kontrolle des sicheren Schlusses!
     \end{enumerate}
     \item Angurten: \ldots
     \item \ldots
 \end{enumerate}
 }{}{}{}{}
 .
 \jALTende
 
 \jNEU
  \Layout{2}{Vorgehen}{
 \begin{enumerate}
     \item \ldots
     \item \EE{Öffnen} der Trage: \ldots
     \item \EE{Aufschaufeln}:
     \begin{enumerate}
       \item 1. Methode: Fußseite offen, Patient wird von oben
      kommend wie auf eine Schere aufgeladen, die Trage unter ihm
      geschlossen, dabei
      stabilisiert ein 3. Helfer
      den Kopf. Abschließend Kontrolle des sicheren Schlusses!
      
      Diese Methode ist besonders bei Becken- und
      Oberschenkelverletzungen sinnvoll und der zweiten
      Methode vorzuziehen da der Patient nicht am Becken
      bewegt wird.
      \item 2. Methode: Beide Teile sind getrennt.
      \begin{enumerate}
        \item Der Patient wird leicht an
        Schulter und Hüfte abwechselnd zur Seite gedreht, dabei
        stabilisiert ein 3. Helfer
        den Kopf und folgt dabei der Drehbewegung achsengerecht.
        Die jeweiligen Tragenteil werden während der Drehbewegung
        unter den Patienten geschoben.
        \item Dann werden die Hälften zuerst am Kopf-, dann am Fußende
        geschlossen.
        \item Anschließend wird eine Kontrolle des sicheren Schlusses durchgeführt
      \end{enumerate}
     \end{enumerate}
     \item Angurten: \ldots
     \item \ldots
 \end{enumerate}
 }{}{}{}{}
 .
\jNEUende

\section{[VIT]}
 \paragraph{Überblick}
 
 Die Vitalfunktionen ermöglichen die
 Funktion des Körpers. Man unterscheidet grundlegende
 Vitalfunktionen 1. Ordnung (Bewusstsein, Atmung, Kreislauf),
 deren Ausfall binnen kurzem zum Tod führen können, und
 Vitalfunktionen 2. Ordnung, deren Störung längere Zeit
\deleted{toleriert werden kann:}

\added{toleriert werden kann. Die Untersuchung der Vitalfunktionen
1. Ordnung sind fester Bestandteil des Einschätzungsblockes
(\FullPrettyRef{topic:einschaetzungsblock}).}

\subsection{\deleted{Bewusstsein}}
\subsection{Vitalfunktionen 1. Ordnung}

\jNEU
\Layout{0}{{\TxQuerverweise}}{%
Die Untersuchung der Vitalfunktionen
1. Ordnung sind aufgrund ihrer Überlebenswichtigkeit fester
Bestandteil des Einschätzungsblockes und werden ebendort
näher erläutert:
\begin{itemize}
  \item \FullPrettyRef{topic:einschaetzungsblock}
\end{itemize}
}{
\begin{itemize}
  \item \FullPrettyRef{topic:einschaetzungsblock}
\end{itemize}
}{}{}{}
.
\jNEUende 
 
\subsubsection{Bewusstsein}
\jNEU
\Layout{0}{{\TxQuerverweise}}{%
Bewusstsein ist ein Oberbegriff für u.\,a. Wachheit,
Orientierung, Aufmerksamkeit, Auffassungsgabe, Denkverlauf und
Merkfähigkeit \cite{Pschy259}. Es ist für das Überleben des
Organismus essentiell und  stellt somit eine Vitalfunktion
1.
Ordnung dar. Aufgrund der Wichtigkeit des Bewusstseins ist
die Bewusstseinskontrolle fester Bestandteil des
Einschätzungsblockes und wird dort näher ausgeführt:
\begin{itemize}
  \item \FullPrettyRef{topic:bewusstsein}.
\end{itemize}
}{
\begin{itemize}
  \item \FullPrettyRef{topic:bewusstsein}
\end{itemize}
}{}{}{}
.
\jNEUende

\subsection{\deleted{Atmung}}

\subsubsection{Atmung}
 
 \jNEU
\Layout{0}{{\TxQuerverweise}}{%
Die Atmung (und damit auch der sichere Atemweg) ist
überlebenswichtig und stellt somit eine Vitalfunktion 1.
Ordnung dar. Die Kontrolle der Atemwege und der Atmung sind
als Punkte  {\FontKeyboardFamily{A}} und
{\FontKeyboardFamily{B}} fester Bestandteil des
Einschätzungsblockes und werden dort näher ausgeführt:
\begin{itemize}
  \item {\FontKeyboardFamily{A}} Atemweg (Airway):
  \FullPrettyRef{topic:atemweg}
  \item {\FontKeyboardFamily{B}} Atmung (Breathing):
  \FullPrettyRef{topic:atmung}
\end{itemize}
}{
\begin{itemize}
  \item {\FontKeyboardFamily{A}} Atemweg:
  \FullPrettyRef{topic:atemweg}
  \item {\FontKeyboardFamily{B}} Atmung:
  \FullPrettyRef{topic:atmung}
\end{itemize}
}{}{}{}
.
 \jNEUende
 
 \subsection{\deleted{Herz-\,/\,Kreislauf}}

\subsubsection{Kreislauf}
\jNEU
\Layout{0}{{\TxQuerverweise}}{%
Der Rkeislauf ist
überlebenswichtig und stellt somit eine Vitalfunktion 1.
Ordnung dar. Die Kontrolle des Kreislaufs ist als Punkt 
{\FontKeyboardFamily{C}} fester
Bestandteil des Einschätzungsblockes und wird dort näher
ausgeführt:
\begin{itemize}
  \item {\FontKeyboardFamily{C}} Kreislauf (Circulation):
  \FullPrettyRef{topic:kreislauf}
  \item Schock:
  \FullPrettyRef{topic:schock}
\end{itemize}
}{
\begin{itemize}
  \item {\FontKeyboardFamily{C}} Kreislauf:
  \FullPrettyRef{topic:kreislauf}
  \item Schock:
  \FullPrettyRef{topic:schock}
  \end{itemize}
}{}{}{}
.
\jNEUende

\subsection{Wasser-/Elektrolythaushalt}
\Layout{0}{Verteilungsräume}
 {%
 Das Körperwasser macht ca. 60\PC* des Körpergewichtes
 aus und teilt sich auf:
 \begin{itemize}
   \item \T{IntraZellularRaum} (IZR), 40\PC* KG
   \item \T{ExtraZellularRaum} (EZR), 20\PC* KG
   \begin{itemize}
     \item \index{Zwischengewebsraum}Zwischengewebsraum,
     $\sim$\,16\PC*: \enquote{freier}
     Raum zwischen den Zellen im Gewebe
     \item \index{Gefäßraum}Gefäßraum $\sim$\,4\PC*:
     Raum in den Blutgefäßen
   \end{itemize}
 \end{itemize}
 Die einzelnen Verteilungsräume haben Trennungen: Der
 Intrazellularraum (\Abk{IZR}) wird vom Extrazellularraum
 (\Abk{EZR}) durch die Zellmembran getrennt, der
 Zwischengewebsraum vom Gefäßraum durch die Blutgefäßwand.
}{%
\jALT
 \begin{itemize}
    \item IZR -- EZR:    Zellmembran (semipermeable Membran,
    s.\,u.)
    \end{itemize}
    .
    \jALTende
    
    \jNEU
    \begin{itemize}
  \item \EE{I}ntra\EE{Z}ellular\EE{R}aum (IZR), 40\PC* KG
  \item \EE{E}xtra\EE{Z}ellular\EE{R}aum (EZR), 20\PC* KG
  \begin{itemize}
    \item Zwischengewebsraum,
    $\sim$\,16\PC*
    \item Gefäßraum $\sim$\,4\PC*
  \end{itemize}
  \item Begrenzungen:
  \begin{itemize}
    \item IZR -- EZR:    Zellmembran 
     \item Zwischengewebsraum -- Gefäßraum:  Blutgefäßwand
  \end{itemize}
 \end{itemize}
 .
 \jNEUende
  }{}{}{}
  
  \clearpage
  \jALT
 \subsection{Wasser- und Elektrolythaushalt}
\index{Wasserhaushalt}
\index{Volumenhaushalt}
\index{Elektrolythaushalt}
\label{topic:wasser-elektrolythaushalt}

\Layout{0}{Exkurs}
{%
Das Blutvolumen setzt sich zusammen aus dem \E{Plasma},
welches zum Großteil aus Wasser besteht, und den
\E{zellulären Bestandteilen} (Blutkörperchen).
Der Anteil des Blutvolumens beträgt beim Erwachsenen
\VonBis{7}{8}\PC*, beim Neugeborenes \VonBis{8}{9}\PC* des Körpergewichts.
\E{Verluste bis zu ca. 20\PC*} sind prinzipiell ohne
Funktionsbeeinträchtigung tolerierbar, allerdings kommt es
bald zu Schäden in der Endstrombahn.

Osmose und Diffusion sorgen zusammen mit den Elektrolyten
für die richtige Verteilung der Flüssigkeit in den
Verteilungsräumen und damit für die Volumenregulation.
Darüber hinaus steuern
Hormone\ZFootnote{Renin, Angiotensin, Aldosteron} zusammen
mit den Nieren die planmäßige Ausscheidung von Flüssigkeit.
}{
\begin{compactitem}
  \item Blutvolumen = Plasma(wasser) + zelluläre Bestandteile
  \item Erwachsener \VonBis{7}{8}\PC*, Neugeborenes\,/\,Säugling \VonBis{8}{9}\PC* des Körpergewichts
  \item Verluste bis 20\PC* verkraftbar
  \item Osmose und Diffusion verteilen Wasser
  \item Hormone {\&} Nieren sorgen für Ausscheidung
\end{compactitem}
}
{}{}{}

\subsection{Osmose und Diffusion}
\index{Osmose|Z}
\index{Diffusion|Z}
\label{topic:osmose}
\label{topic:diffusion}
 
 \Layout{2}
{\TxBeschreibung}
{\DefinitionInPara{Unter Osmose versteht man den Fluss
 von Molekülen durch eine halbdurchlässige Membranen.} 
 Im menschlichen Körper spielt die Osmose eine wesentliche
Rolle in der Flüssigkeitsverteilung.}
{}
{}
{}
{}


\AuthNotes{GAB: Osmose und Diffusion: Grundsätzlich wichtig
und relevant, allenfalls bessere Erklärung, aber bleibt in
Normalschrift, nicht heller. [GAB]
 
KOCH: Sollte alles grau werden}

\Layout{2}{Flüssigkeitsverschiebung durch Diffusion}
{%
%
 Geladene Teilchen können die Zellmembran nicht selbstständig
 durchdringen, Wasser dagegen schon. Man nennt die
 Zellmembran daher eine \E{semipermeable
halbdurchlässige Membran}\ZFootnote{\I{semipermeable
Membran}\index{Membran!semipermeable}}.
 Die Natur möchte, dass in einer Flüssigkeit überall \E{die
 gleiche Konzentration geladener Teilchen} herrscht. Da
 die Zellmembran (und teilweise auch die Blutgefäßwand) nur
 das Wasser durchlässt, nicht jedoch die geladenen Teilchen,
verschiebt sich nur das Wasser von einem Raum in den
anderen. Man kann
auch sagen:

\Fazit{\enquote{Das Wasser folgt dem Konzentrationsgefälle.}}

\Z{Man kann es sich so vorstellen, dass die Elektrolyte einen
\enquote{Druck} ausüben und das Wasser dazu bringen \enquote{die Seiten
 zu wechseln}. Man nennt diesen Druck den
\enquote{\index{osmotischer Druck}\T{osmotischen Druck}}.} 
\Z{Daneben gibt es den \E{onkotischen} Druck, welcher durch
Eiweißbestandteile ausgeübt wird.\ZFootnote{Der durch
 Eiweißbestandteile hervorgerufene onkotische Druck wird bei
 speziellen \enquote{kolloidalen} Infusionslösungen ausgenützt. Sie
 enthalten große Eiweiße, welche die Blutgefäßwand
 normalerweise nicht passieren können und folglich Wasser vom
 Zwischengewebsraum in den Blutgefäßraum \enquote{ziehen}. Der
 verbreitetste Wirkstoff ist die Hydroxyethylstärke (HAES),
 Beispiele für damit erzeugte Infusionslösungen sind
 Elohaest{\texttrademark}, Voluven{\texttrademark} und
HyperHAES{\texttrademark}.}}
}{
 
 }
% {./images/osmose.png}
 {}
 {}
 {}
 
  
 \begin{PictureSeriesWide}{}{Bilderserie: \EEE{Flüssigkeit im Körper}}
 \PictureSeriesRowWide{2}
 {./images/wikipedia-pd/Semipermeable_membrane.png}
 {Gefäßraum und Zwischengewebsraum, getrennt durch die
 Blutgefäßwand. Es schwimmt viel herum: Die grünen, gelben
 blauen und violetten Kugeln stellen Ionen und Moleküle dar.
 Die roten Scheiben sind rote Blutkörperchen. 
 \SourcePicture{WMC/pidalka44}}
 {./images/wikipedia-pd/Osmotische_druk.png}
 {Diffusion \SourcePicture{WMC (Sander van der Molen)/PD}}
 {}
 {}
 {}
 {}
 \end{PictureSeriesWide}
 

\Z{Wie kommen die Elektrolyte in die einzelnen Räume wenn
doch nur Wasser die Wand passieren kann? Dafür gibt es in
der Zelle spezielle \enquote{Pumpen}
 \Z{(Transportproteine, z.\,B. die
 \I{Natrium-Kalium-ATPase})} die die Ionen zwischen den
 intra- und extrazelluärraum befördern.
 }
 
Die Steuerung der richtigen Verteilung der verschiedenen
Ionen ist recht kompliziert aber lebenswichtig. Eine
falsche Verteilung der Elektrolyte kann schlimme Folgen
haben (zB. Herzrhythmusstörungen).
 
 \Fazit{Die richtige Verteilung der Elektrolyte im Körper
 ist lebenswichtig. Elektrolytstörungen können
 lebensgefährlich sein. }
 
 \Layout{0}{Häufige Ursachen von
 Elektrolytstörungen}
 {%
 \index{Elektrolytstörungen}
 \label{topic:elektrolytstoerungen}
 \label{topic:dialyse-elektrolyte}
 %
 \begin{itemize}
     \item Niereninsuffizienz (gestörte Ausscheidung)
     \begin{itemize}
         \item Besonders bei
         Dialysepatienten\index{Dialyse!Elektrolytstörungen}
     \end{itemize}
     \item Durchfall, Erbrechen (Verlust von Elektrolyten)
     \item Medikamente
     \begin{itemize}
       \item \I[Insulin!Elektrolytstörungen]{Insulin}: Bei
       der Behandlung einer Hyperglykämie (zu hoher Blutzuckerspiegel) 
       mittels Insulin kann es
       zu lebensbedrohlichen Elektrolytstörungen kommen!
     \end{itemize}
     \item Vergiftung
     %\item \Z{Hormonentgleisungen}
 \end{itemize}
 }{%
 \begin{itemize}
     \item Niereninsuffizienz, Dialysepatienten
     \item Durchfall, Erbrechen (Verlust von Elektrolyten)
     \item Medikamente (Insulin!), Vergiftung
     %\item \Z{Hormonentgleisungen}
 \end{itemize}
 }{}{}{}
.
 \jALTende
 
 \clearpage
 \jNEU
 \Layout{2}
{Osmose und Diffusion}
{%
\index{Osmose|Z}%
\index{Diffusion|Z}%
\label{topic:osmose}%
\label{topic:diffusion}%
\DefinitionInPara{Unter Osmose versteht man den Fluss
 von Molekülen durch eine halbdurchlässige Membranen.} 
 Im menschlichen Körper spielt die Osmose eine wesentliche
Rolle in der Flüssigkeitsverteilung.
 
 Geladene Teilchen können die Zellmembran nicht selbstständig
 durchdringen, Wasser dagegen schon. Man nennt die
 Zellmembran daher eine \E{semipermeable
(halbdurchlässige) Membran}\index{Membran!semipermeable}.
 Die Natur möchte, dass in einer Flüssigkeit überall \E{die
 gleiche Konzentration geladener Teilchen} herrscht. Da
 die Zellmembran (und teilweise auch die Blutgefäßwand) nur
 das Wasser durchlässt, nicht jedoch die geladenen Teilchen,
kann dies nur erreicht werden, indem Wasser vom Raum der
niedrigeren Konzentration in den mit einer höheren
Konzentration übertritt.%
\ZFootnote{Man kann es sich so vorstellen, dass die
Elektrolyte einen \enquote{Druck} ausüben und das 
Wasser dazu bringen \enquote{die Seiten
 zu wechseln}. Man nennt diesen Druck den
\enquote{\index{osmotischer Druck}\T{osmotischen Druck}}.
Daneben gibt es den \E{onkotischen} Druck, welcher durch
Eiweißbestandteile ausgeübt wird. Der durch
 Eiweißbestandteile hervorgerufene onkotische Druck wird bei
 speziellen \enquote{kolloidalen} Infusionslösungen ausgenützt. Sie
 enthalten große Eiweiße, welche die Blutgefäßwand
 normalerweise nicht passieren können und folglich Wasser vom
 Zwischengewebsraum in den Blutgefäßraum \enquote{ziehen}. Der
 verbreitetste Wirkstoff ist die Hydroxyethylstärke (HAES),
 Beispiele für damit erzeugte Infusionslösungen sind
 Elohaest{\texttrademark}, Voluven{\texttrademark} und
HyperHAES{\texttrademark}.} 
Man kann
auch sagen:

\Fazit{\enquote{Das Wasser folgt dem Konzentrationsgefälle.}}
 
 }
{
}
 {}
 {}
 {}
 
 
 
\AuthNotes{GAB: Osmose und Diffusion: Grundsätzlich wichtig
und relevant, allenfalls bessere Erklärung, aber bleibt in
Normalschrift, nicht heller. [GAB]

KOCH: Sollte alles grau werden}

 
 \begin{PictureSeriesWide}{}{Bilderserie: \EEE{Flüssigkeit im Körper}}
 \PictureSeriesRowWide{2}
 {./images/wikipedia-pd/Semipermeable_membrane.png}
 {Gefäßraum und Zwischengewebsraum, getrennt durch die
 Blutgefäßwand. Es schwimmt viel herum: Die grünen, gelben
 blauen und violetten Kugeln stellen Ionen und Moleküle dar.
 Die roten Scheiben sind rote Blutkörperchen. 
 \SourcePicture{WMC/pidalka44}}
 {./images/wikipedia-pd/Osmotische_druk.png}
 {Diffusion \SourcePicture{WMC (Sander van der Molen)/PD}}
 {}
 {}
 {}
 {}
 \end{PictureSeriesWide}
 
\Layout{2}
{Verteilung der Elektrolyte}
{%
Die Steuerung der richtigen Verteilung der
Elektrolyte zwsichen den Verteilungsräumen ist recht
kompliziert.%
\ZFootnote{Verteilung der Elektrolyte:  
Dafür gibt es in der Zelle spezielle \enquote{Pumpen}
 \Z{(Transportproteine, z.\,B. die
 \I{Natrium-Kalium-ATPase})} die die Ionen zwischen den
 intra- und extrazelluärraum befördern.
}% 
Eine
falsche Verteilung der Elektrolyte kann schlimme Folgen
haben (z.\,B. Herzrhythmusstörungen). Zu solchen
Verteilungsstörungen kann es z.\,B. nach einer Blutwäsche
(Dialyse) oder bei Vergiftungen kommen.
 }
{}
{}
{}
{}
 
 \Fazit{Die richtige Verteilung der Elektrolyte im Körper
 ist lebenswichtig. Elektrolytstörungen können
 lebensgefährlich sein. }
 
 \Layout{0}{Häufige Ursachen von
 Elektrolytstörungen}
 {%
 \index{Elektrolytstörungen}
 \label{topic:elektrolytstoerungen}
 \label{topic:dialyse-elektrolyte}
 %
 \begin{itemize}
     \item Niereninsuffizienz (gestörte Ausscheidung)
     \begin{itemize}
         \item Besonders bei
         Dialysepatienten\index{Dialyse!Elektrolytstörungen}
     \end{itemize}
     \item Durchfall, Erbrechen (Verlust von Elektrolyten)
     \item Medikamente
     \begin{itemize}
       \item \I[Insulin!Elektrolytstörungen]{Insulin}: Bei
       der Behandlung einer Hyperglykämie (zu hoher Blutzuckerspiegel) 
       mittels Insulin kann es
       zu lebensbedrohlichen Elektrolytstörungen kommen!
     \end{itemize}
     \item Vergiftung
     %\item \Z{Hormonentgleisungen}
 \end{itemize}
 }{%
 \begin{itemize}
     \item Niereninsuffizienz, Dialysepatienten
     \item Durchfall, Erbrechen (Verlust von Elektrolyten)
     \item Medikamente (Insulin!), Vergiftung
     %\item \Z{Hormonentgleisungen}
 \end{itemize}
 }{}{}{}
 
\Layout{0}
{Ausscheidung und Flüssigkeitsverlust}
{%
Der Körper kann Flüssigkeit (und auch Elektrolyte) entweder
durch (gesteuerte) Ausscheidung (Volumenregulation) oder
durch Flüssigkeitsverlust verlieren:
 
\begin{itemize}
  \item \EE{Ausscheidung}: Über die \EE{Nieren}, geregelt
  u.\,a. durch Hormone\ZFootnote{Renin, Angiotensin, Aldosteron}
  \item \EE{Flüssigkeitsverlust}: Z.\,B.:
  \begin{itemize}
    \item Stuhl
    \item Haut, bzw. Schweiß
    \item Atmung: Der Flüssigkeitsverlust bei der Atmung
    wird häufig unterschätzt. Fiebernde Patienten können bis
    zu 2 Liter an Flüssigkeit nur über die Atmung verlieren!
  \end{itemize}
\end{itemize}
}
{%
\begin{itemize}
  \item \EE{Ausscheidung}: {Nieren}
  \item \EE{Flüssigkeitsverlust}: Z.\,B.:
  \begin{itemize}
    \item Stuhl
    \item Haut, Schweiß
    \item Atmung!
  \end{itemize}
\end{itemize}
}
{}
{}
{}
 .
 \jNEUende

\subsection{Säure-Basen-Haushalt}
\jALT
\subsection{Allgemeines}

\Layout{abstract}{}
{Der Körper muss einen ausgewogenen Säure-Basenhaushalt
(pH-Wert) haben. Puffersysteme sind wichtige Intrumente zur
Regulation. Der pH-Wert hängt eng mit der Atmung zusammen:
Säure kann zu Wasser und \ce{CO2} abgebaut werden. Bei zu
geringer Atmung sammelt sich im Körper Säure, bei
viel Atmung wird mehr Säure verbraucht. Umgekehrt erhöht
sich bei Übersäuerung der Atemantrieb. 

}{}{}{}{}
.
\jALTende

\jALT 
\Layout{0}{pH-Wert}{
 \index{pH-Wert}
 Der pH-Wert gibt an, wie \E{sauer} oder \E{basisch} eine
Lösung ist. Ein pH-Wert {\textless}\,7 gilt als \E{sauer},
 ein pH\,=\,7 als \E{neutral} und ein ph\,{\textgreater}\,7
als \E{basisch}.

Der Körper ist da
sehr empfindlich, der pH-Wert muss sich in engen Grenzen
 bewegen. Schon geringe Abweichungen können große Folgen
 haben. Bei einer Übersäuerung sprechen wir von einer
\T{Azidose}, wenn der Patient basisch wird,
von einer \T{Alkalose}.
 
% \Fazit{Bei einer \EE{Übersäuerung} sprechen wir von einer
% \T{Azidose}, wenn der Patient \EE{basisch} wird von einer
% \T{Alkalose}.}

\Z{Der richtige pH-Wert ist lebenswichtig. Er wird beim
 Menschen normalerweise vom Körper sehr konstant um ca 7,4
 gehalten. Bereits einen pH-Wert unter 7,35 (!)  bezeichnet
 man als \E{Azidose}, bzw. über 7,45  als \E{Alkalose}.
 Leider gibt es präklinisch normalerweise keine Möglichkeit
 den pH-Wert zu messen, jedoch hat die Atmung bzw. Beatmung
 sowie Erkrankungen
 großen Einfluss auf ihn, siehe unten.}
 }{
 \begin{itemize}
  \item Gibt an, wie sauer etwas ist.
  \item pH\,{\textless}\,7 {\dotfill} sauer, 
  pH\,=\,7 ist neutral,
  \item pH\,{\textgreater}\,7 {\dotfill} basisch
 \end{itemize}
 }{}{}{}
.
\jALTende

\jNEU
\Layout{0}{pH-Wert}
{
 \index{pH-Wert}
 Der pH-Wert gibt an, wie \E{sauer} oder \E{basisch} eine
Lösung ist.\ZFootnote{Ein pH-Wert {\textless}\,7 gilt als \E{sauer},
 ein pH\,=\,7 als \E{neutral} und ein ph\,{\textgreater}\,7
als \E{basisch}.} 
Der pH-Wert im Blut muss sich in engen Grenzen
 bewegen. Schon geringe Abweichungen können große Folgen
 haben. Bei einer Übersäuerung sprechen wir von einer
\T{Azidose}, wenn der Patient basisch wird
von einer \T{Alkalose}.%
\ZFootnote{Der richtige pH-Wert ist lebenswichtig. Er wird beim
 Menschen normalerweise vom Körper sehr konstant um ca 7,4
 gehalten. Bereits einen pH-Wert unter 7,35 (!)  bezeichnet
 man als \E{Azidose}, bzw. über 7,45  als \E{Alkalose}.
 Leider gibt es präklinisch normalerweise keine Möglichkeit
 den pH-Wert zu messen, jedoch hat die Atmung bzw. Beatmung
 sowie Erkrankungen
 großen Einfluss auf ihn.}
 }{
 \begin{itemize}
  \item Gibt an, wie sauer oder basisch ein Stoff ist.
  \item Übersäuerung: Azidose
  \item Zu basisch: Alkalose
%   \item pH\,{\textless}\,7 {\dotfill} sauer, 
%   pH\,=\,7 ist neutral,
%   \item pH\,{\textgreater}\,7 {\dotfill} basisch
 \end{itemize}
 }{}{}{}
.
\jNEUende

\jALT
 \subsubsection*{Zusammenfassung}
 
 
 \Fazit{Die \EE{richtige Verteilung} von Wasser im Körper ist
lebenswichtig.}

\Fazit{Mittels Osmose und Diffusion verteilt der Körper das
Wasser.} 

\Fazit{Die \EE{Elektrolyte} haben einen großen Einfluss auf
die
  Wasserverteilung: Nur wenn die Elektrolyte richtig
  verteilt sind, kann das Wasser richtig verteilt werden.}
\Fazit{Elektrolyte haben auch
  eine elektrische Funktion.}
\Fazit{Ein ausgeglichener Säure-Basen-Haushalt ist wichtig, damit chemische
  Abläufe im Körper richtig funktionieren können.}
  
 \Fazit{Besonders die \EE{Atmung} und der \EE{Stoffwechsel} können dem
   Säure-Basen-Haushalt Streiche spielen: \newline
   Wer zu wenig atmet wird sauer. \newline
   Wer zu viel atmet wird basisch und spürt ein Kribbeln
 in den Fingern. \newline
 Wer Probleme mit dem Stoffwechsel hat kann auch
     sauer oder basisch werden.}
\jALTende

\jNEU
 \subsubsection*{Zusammenfassung}
 
 
 \Fazit{Die \EE{richtige Verteilung} von Wasser im Körper ist
lebenswichtig. Mittels Osmose und Diffusion verteilt der
Körper das Wasser. Die \EE{Elektrolyte} haben dabei einen großen
Einfluss: Nur wenn die Elektrolyte richtig
  verteilt sind, kann das Wasser richtig verteilt werden.
  Daneben haben sie auch eine elektrische
  Funktion.}
\Fazit{Der Körper muss einen ausgewogenen
Säure-Basenhaushalt (pH-Wert) haben, damit chemische
  Abläufe im Körper richtig funktionieren können.. 
  Puffersysteme sind wichtige Instrumente zur
Regulation. Der pH-Wert hängt eng mit der Atmung zusammen:
Säure kann zu Wasser und \ce{CO2} abgebaut werden. Bei zu
geringer Atmung sammelt sich im Körper Säure, bei
viel Atmung wird mehr Säure verbraucht. Umgekehrt erhöht
sich bei Übersäuerung der Atemantrieb. }
 \Fazit{Besonders die \EE{Atmung} und der \EE{Stoffwechsel} können dem
   Säure-Basen-Haushalt Streiche spielen: \newline
   Wer zu wenig atmet wird sauer. \newline
   Wer zu viel atmet wird basisch und spürt ein Kribbeln
 in den Fingern. \newline
 Wer Probleme mit dem Stoffwechsel hat kann auch
     sauer oder basisch werden.}
     .
\jNEUende


SO LONG.......... AND THANKS FOR ALL THE FISH


\chapter{Änderungen im Maßnahmenkatalog}

Keine Änderungen.

\chapter{Neue Literatur}
\begin{itemize}
  \item \fullcite{Thierbach2004}%~~~\bibentry{Thierbach2004}
\item \fullcite{Barnes2001}%~~~\bibentry{Barnes2001}
\item \fullcite{Mann2013}%~~~\bibentry{Mann2013}
\item \fullcite{Schalk2012}%~~~\bibentry{Schalk2012}
\item \fullcite{Schalk2012a}%~~~\bibentry{Schalk2012a}
\item \fullcite{Bernhard2012}%~~~\bibentry{Bernhard2012}
\item \fullcite{Cook2011}%~~~\bibentry{Cook2011}
\item \fullcite{Wang2011}%~~~\bibentry{Wang2011}
\item \fullcite{Dengler2011}%~~~\bibentry{Dengler2011}
\item \fullcite{Schalk2010}%~~~\bibentry{Schalk2010}
\item \fullcite{Cobas2009}%~~~\bibentry{Cobas2009}
\item \fullcite{Wiese2009}%~~~\bibentry{Wiese2009}
\item \fullcite{Jones2004}%~~~\bibentry{Jones2004}
\item \fullcite{Grmec2002}%~~~\bibentry{Grmec2002}
\item \fullcite{Lyon2010}%~~~\bibentry{Lyon2010}
\item \fullcite{Pelucio1997}%~~~\bibentry{Pelucio1997}
\item \fullcite{Jemmett2003}%~~~\bibentry{Jemmett2003}
\item \fullcite{Katz2001}%~~~\bibentry{Katz2001}
\item \fullcite{Komatsu2010}%~~~\bibentry{Komatsu2010}
\item \fullcite{Mulcaster2003}%~~~\bibentry{Mulcaster2003}
\item \fullcite{Konrad1998}%~~~\bibentry{Konrad1998}
\item \fullcite{Timmermann2012}%~~~\bibentry{Timmermann2012}
\end{itemize}


\chapter{Diskussionen}
\section{Inhaltlich}


\section{Formal/Layout}


\section{Reviewprotokoll}


\chapter{Ergebnis/TO DO}
\begin{enumerate}
  \item Review des \CC{[PAM]}-Kapitels im Zuge des 44. Reviews Ende
  Oktober/Anfang November
  \item Kapitel \CC{[ASS]}: Durchstechampullen erwähnen
\end{enumerate}

\appendix
\chapter{Anhang}
Tischvorlage Probekapitel \CC{[PAM]}

\MultiColListEnd{toc}

%\bibliography{BiblioDB-Med-Misc}
\printbibliography
\end{document}